\documentclass[preprint,aps,nofootinbib,a4paper,superscriptaddress,longbibliography,amsfonts, amssymb,amsmath,titlepage]{revtex4-1}
\ProvidesPackage{jheppub}[2013/03/21 r534]
\numberwithin{equation}{section} % JHEP style eqn
%\input{CQGformat.input}
%s\documentclass[cqg]{iopart}
\usepackage{amssymb}
%\pdfoutput=1
%\usepackage{iopams, setstack}
\usepackage{bm}
\usepackage{bbm}
\usepackage{color}
\usepackage{microtype}
\usepackage{IEEEtrantools}
\usepackage{comment}
\usepackage{mathrsfs}
\usepackage{amsmath}
\usepackage{amsthm}
\usepackage{enumitem}
\usepackage[normalem]{ulem}
%\usepackage[numbers,square,sort&compress]{natbib}
\usepackage[pdftex]{hyperref}
%\bibliographystyle{iopart-num}
\usepackage[T1]{fontenc}
% Some DOIs in MyReferences.bib contain underscores (e.g. Pound:2021qin),
% which apsrev4-1 typesets in text mode and which therefore break the build.
% The underscore package makes a bare _ printable outside maths.
\usepackage[strings]{underscore}

\newcommand{\PZ}[1]{\textcolor{red}{\textsf{[PZ: #1]}}}
\newcommand{\newtext}[1]{\textcolor{cyan}{#1}}

\def\half{\tfrac{1}{2}}
\def\pihalf{\tfrac{\pi}{2}}


\def\pheq{\phantom{=}}
\newcommand{\preindex}[2]{\,{}_{#2}#1}
\usepackage[toc,title,page,titletoc]{appendix}
%  vectors and diff Forms
\def\bF{\boldsymbol{F}}
\def\bd{\boldsymbol{\mathrm{d}}}
\newcommand{\x}[1]{\xi_{#1}^{\circ}} % gauge vector

% GHP weight
\def\GHPwt{\ensuremath{\circeq}}
% NP SpinCoeffs
\def\kap{\kappa'}
\def\sip{\sigma'}
\def\rhp{\rho'}
\def\rhph{{\rho'}^\circ}
\def\tap{\tau'}
\def\tabp{\bar{\tau}'}
\def\bep{\beta'}
\def\epp{\epsilon'}
\def\rhp{\rho'}
\def\sib{\bar\sigma}
\def\epb{\bar\epsilon}
\def\kab{\bar\kappa}
\def\beb{\bar\beta}
\def\rhb{\bar\rho}
\def\pib{\bar\pi}
\def\alb{\bar\alpha}
\def\gab{\bar\gamma}
\def\mub{\bar\mu}
\def\lab{\bar\lambda}
\def\nub{\bar\nu}
\def\tab{\bar\tau}
% NP leg
\def\lmu{l^\mu}
\def\nmu{n^\mu}
\def\mmu{m^\mu}   % PZ: was n^\mu in main.tex -- typo, fixed here
\def\mbmu{{\bar m}^\mu}
\def\mb{\bar m}
%
% Held scalars
%
\renewcommand{\H}[1]{{#1}^{\circ}}
\newcommand{\Held}[1]{{#1}^{\circ}}
\def\tah{{\tau}^\circ}
\def\tabh{{\bar\tau}^\circ}
\def\tauH{\tau^{\circ}}
\def\taubH{\bar{\tau}^{\circ}}
\def\Psih{ {\Psi}^\circ}
\def\Psibh{ {\bar{\Psi}}^\circ}
\def\PsiH{\Psi^{\circ}}
\def\PsibH{\bar\Psi^{\circ}}
\def\rhoph{{\rho}^{\prime \circ} }
\def\rhopH{\rho^{\prime \circ}}
\def\rhopbh{{\bar\rho}^{\prime \circ} }
\def\tauHsq{\tau^{\circ}{}^{2}}
\def\Oh{{\Omega}^\circ}

% Nicer looking GHP Ops FROM GHZ PAPER
\font\ec=ecrm0800 at 12pt
\def\thpH{\tilde{\th}'}
\def\edthH{\tilde\edth}
\def\edthpH{\tilde\edth'}
\def\th{\hbox{\ec\char'336}}
\def\edth{\hbox{\ec\char'360}}
\def\thp{\hbox{\ec\char'336}'}
\def\edthp{\hbox{\ec\char'360}'}
\def\Dbar{\tilde \edth} % Held edth
\def\Nbar{\tilde \th'}
\def\Edth{\hat{\edth}} % raising operator
\def\EdthBar{\bar{\Edth}} % lowering operator on sYlm
%
%  SWS Harmonics
%
\newcommand{\C}[1]{\lambda_{l,#1}}
\newcommand{\thA}[1]{\theta^{A}_{#1}}
\def\thA{\ensuremath{\theta,\varphi}}
\def\Omp{\Omega_\p}
\newcommand{\CGCsqrt}[3]{\sqrt{\frac{(2{#1}+1)(2{#2}+1)(2{#3}+1)}{4\pi}}}
\newcommand{\wthreej}[6]{%
  \ensuremath{%
  \begin{pmatrix}
      #1 & #2 & #3 \\
      #4 & #5 & #6
    \end{pmatrix}
  }%
}
% mode indices
\newcommand{\emm}{m}
\newcommand{\ess}{s}
\def\wlm{\ell \emm \omega}
% spin weighted spherical harmonics
\newcommand{\Y}[1]{{}_{#1}Y_{\ell, m}}
\newcommand{\bY}[1]{{}_{#1}\bar{Y}_{\ell, m}}
%\newcommand{\Ym}[2]{{}_{#1}Y_{l{#2}}}
\newcommand{\Ym}[2]{\, {}_{#1}Y_{\ell, #2}}
\newcommand{\Ylm}[3]{\, {}_{#1}Y_{#2, #3}}
\newcommand{\bYlm}[3]{\, {}_{#1} \bar{Y}_{#2, #3}}

% Sph. harmonic defs (Peter)
\def\l{\ensuremath{\ell}}
\def\lm{{\ensuremath{\ell \emm}}}
% at the particle
\newcommand{\Yp}[1]{\, {}_{#1}Y^\textrm{p}_{\ell ,\emm}}
\newcommand{\Yplm}[3]{\, {}_{#1}Y^\textrm{p}_{#2 ,#3}}
\newcommand{\bYp}[1]{\, {}_{#1}\bar{Y}^\textrm{p}_{\ell, \emm}}
\def\YpL{ Y_{\ell}^{\p}}
\def\bYpL{\bar{Y}_{\ell}^{\p}}
\newcommand{\YpsL}[1]{{}_{#1} Y_{\ell}^{\p}}
\newcommand{\bYsL}[1]{ {}_{#1}\bar Y_{\ell}^\p}
%\newcommand{\C}[1]{\lambda_{\ell, #1}}
\newcommand{\Ltwo}[2]{\langle {#1} \vert {#2} \rangle_{L^2} }
\def\iStwo{\int_{\mathbb S^2}}
\def\dOm{\dd \Omega}

\newcommand{\xlm}[1]{\langle\xi_{#1}^{\circ}\rangle_{\ell \emm}} % gauge vector modes

% Helicity basis vectors
\def\hmp{{\hat{m}}^+}
\def\hmm{{\hat{m}}^-}
\def\nUL{{\hat n}^{\langle L \rangle}}
\def\uLL{{\hat n}_{\langle L \rangle}}
\def\hz{{\hat z}}

\def\bk{\mathbf{k}}
\def\bsk{\boldsymbol{k}}
%
% Operators
%

% Chandrasekhar ops
\newcommand{\chandR}[2]{\,{}_{#1} \mathcal{L}^\dagger_{#2}{}}
\newcommand{\chandL}[2]{\,{}_{#1} \mathcal{L}^\dagger_{#2}{}}
% Teukolsky/Wald/Chrzanowski operators
\renewcommand{\S}{{\mathcal S}}
\newcommand{\E}{{\mathcal E}}
\newcommand{\T}{{\mathcal T}}
% Lie derivatives
\newcommand{\Lie}{{\mathcal L}}
\newcommand{\Lcsdag}[2]{ {}_{-{#1}}{\mathcal L}^{\dagger}_{#2}}
\newcommand{\Lcs}[2]{ {}_{{#1}}\mathcal L_{#2}}
%  derivative operators
\def\pd{\partial}
\def\dd{{\rm d}}
\newcommand{\dbd}[1]{\frac{\pd}{\pd #1}}
\renewcommand{\Re}{\operatorname{Re}}
\renewcommand{\Im}{\operatorname{Im}}
\def\pd{\partial}
\newcommand{\abs}[1]{\vert #1 \vert}
%
% Eqns.
%
\def\beq{\begin{equation}}
\def\eeq{\end{equation}}
\def\beqs{\begin{subequations}}
\def\eeqs{\end{subequations}}
\def\bal{\begin{align}}
\def\eal{\end{align}}
\def\non{\nonumber \\}
\def\beq{\begin{equation}}
\def\eeq{\end{equation}}
\def\ud{\mathrm{d}}
\def\beqs{\begin{subequations}}
\def\eeqs{\end{subequations}}


\def\half{\tfrac{1}{2}}

% indices
\def\mn{\ensuremath{{\mu\nu}}}
\def\a{\alpha}
\def\b{\beta}
\newcommand{\stf}[1]{\langle {#1}\rangle}

% coordinates
\def\vphi{\ensuremath{{\varphi}}}
\def\sz{\ensuremath{{\sqrt{1-z^2}}}}
\def\z{\zeta}
\def\zb{\bar\zeta}
\def\e{\eta}
\def\eb{\bar\eta}

% at the particle
\def\p{{\rm p}}
\def\rhop{\rho_{\p}}
\def\rhobp{\bar{\rho}_{\p}}
\def\rp{r_{\p}}
\def\dOm{\ensuremath{\delta^2(\Omega-\Omega_\p)}}
\def\zp{z_\p}
\def\rhbp{\bar{\rho}_\p}

\def\rhom{\rho_{\rm min}}
\def\Tld{T^{(\delta,-)}}
\def\Tldp{T^{(\delta',-)}}
\def\Trd{T^{(\delta,+)}}
\def\Trdp{T^{(\delta',+)}}

\def\Dr{\Delta r}
\def\rmx{r_{\mathrm{max}}}
\def\rmn{r_{\mathrm{min}}}


% Spacing
\newcommand{\nls}{\phantom{\frac{1}{\rho}}}
\newcommand{\sforall}{\quad \forall}

% =====================================================================
%  NEW MACROS ADDED FOR THIS DRAFT
% =====================================================================
% angular eigenvalue combinations
\newcommand{\Bco}[1]{{}_{#1}B_{\l}}                 % Schwarzschild  {}_sB_l
\newcommand{\BcoK}[3]{{}_{#1}B_{#2}(#3)}            % Kerr  {}_sB_{lm}(a w)
\def\LamL{\Lambda_{\l}}                             % (l-1)l(l+1)(l+2)
\def\Dlw{\mathcal{D}_{\l\emm\omega}}                % denominator 2B 2Bbar + 9 M^2 w^2
\def\Cts{\mathfrak{C}}                              % complex Teukolsky-Starobinsky constant
% Held integration constants and seed
\def\fHo{f^{\circ}}
\def\fbHo{\bar f^{\circ}}
\newcommand{\PhiN}[1]{\Phi^{#1\circ}}
\def\Phiadj{\Phi^{\rm adj}}
\def\eHo{e^{\circ}}
\def\dHo{d^{\circ}}
% asymptotic Weyl amplitudes
% Single notation for the Teukolsky boundary amplitudes:
%   \Zinf{s}{modes}  = amplitude at scri   (the Toolkit's "Amplitudes"["I"])
%   \Zhor{s}{modes}  = amplitude at H^-    (the Toolkit's "Amplitudes"["H"])
% The spin weight sits in the left subscript, the boundary in the superscript.
% This replaces the C^{+}/C^{-} notation of the notebooks; see Sec. III A.
\newcommand{\Zinf}[2]{{}_{#1}Z^{\infty}_{#2}}
\newcommand{\Zhor}[2]{{}_{#1}Z^{\rm H}_{#2}}
\def\psifive{\psi^{5\circ}_{0}}
\def\psifourone{\psi^{1\circ}_{4}}
\newcommand{\psifiveM}[1]{\psi^{5\circ}_{0,#1}}       % with mode labels
\newcommand{\psifouroneM}[1]{\psi^{1\circ}_{4,#1}}
% residual gauge vector components
\def\zetal{\zeta^{\circ}_{l}}
\newcommand{\zetalM}[1]{\zeta^{\circ}_{l,#1}}
\def\zetan{\zeta^{\circ}_{n}}
\newcommand{\zetanM}[1]{\zeta^{\circ}_{n,#1}}
\def\zetam{\zeta^{\circ}_{m}}
\newcommand{\zetamM}[1]{\zeta^{\circ}_{m,#1}}
\def\zetamb{\bar\zeta^{\circ}_{m}}
% mode labels
\def\lmnk{\l \emm n k}
\def\wmnk{\omega_{\emm n k}}
% harmonics with explicit spin
\newcommand{\sYY}[3]{{}_{#1}Y_{#2,#3}}
\newcommand{\sSS}[3]{{}_{#1}S_{#2,#3}}
% scri
\def\scri{\mathscr{I}^{+}}
% "equals by definition"
\def\defeq{\mathrel{\mathop{:}}=}

\begin{document}

\title{Asymptotically flat metric reconstruction:\\ a mode-based scheme for gravitational self-force calculations}

\author{Vahid Toomani}
\email{vahid.toomani.2@uni-leipzig.de}
\author{Stefan Hollands}
\email{stefan.hollands@uni-leipzig.de}
\author{Peter Zimmerman}
\email{zimmerator@protonmail.com}

\begin{abstract}
% ---------------------------------------------------------------------
% ROADMAP (abstract).  Three sentences of set-up, one of method, one of
% result.  Keep the emphasis on *practicality*: the only new input needed
% is psi_0^{5 o}, which existing Teukolsky codes already produce.
% ---------------------------------------------------------------------
Metric perturbations obtained from a Hertz potential in an ingoing radiation
gauge generically grow linearly in $r$ at future null infinity, and are
therefore not asymptotically flat in the Bondi sense.  Hollands and Toomani
showed that this growth is pure gauge and can be removed, within the
Green--Hollands--Zimmerman (GHZ) framework, by a residual gauge vector
together with a homogeneous, Held-covariantly constant completion of the
Hertz potential.  In this work we turn that formal construction into a
practical, mode-based algorithm.  We show that the entire completion is
generated by a single Held scalar $\fHo$, obtained algebraically from the
leading asymptotic Weyl amplitude $\psifive$ through a Teukolsky--Starobinsky
inversion, and we give explicit $(\l,\emm,\omega)$ formulas for the four
Held integration constants $\PhiN{0},\dots,\PhiN{3}$ and for the residual
gauge vector $\zeta^a$.  In Schwarzschild the construction is diagonal in
spin-weighted spherical harmonics and every completion field is available in
closed form; in Kerr we give the corresponding $(\l,\emm,n,k)$ scheme, in
which the only additional ingredients are the spheroidal-to-spherical mixing
coefficients and a tridiagonal $\cos\theta$ coupling matrix.  We implement
the scheme numerically and demonstrate the cancellation of the growing
radiation-gauge terms and the recovery of Bondi falloff.
\PZ{Add one closing sentence quoting the actual numerical result once the
runs are in, e.g.\ ``the $O(r)$ coefficients cancel to $10^{-N}$ relative
precision for $\l\le\l_{\max}$''.}
\end{abstract}

\maketitle

\tableofcontents


% =====================================================================
\section{Introduction}
\label{sec:introduction}
% =====================================================================
% ---------------------------------------------------------------------
% ROADMAP.  Four subsections:
%   1.1  why we need the metric (2nd order SF, LISA)
%   1.2  what goes wrong with CCK-Ori (the O(r) problem)
%   1.3  what Hollands-Toomani fixed, and what is still missing (a mode
%        algorithm) -- this is the "gap" paragraph
%   1.4  conventions + outline
% Sections 1.1-1.3 below are your existing text, lightly completed where
% sentences were truncated.  Places where I finished a sentence are marked
% with \newtext{...} so you can see exactly what was added.
% ---------------------------------------------------------------------

\subsection{Motivation}\label{sec:motiv}
Extreme-mass-ratio inspirals are among the principal targets of the
Laser Interferometer Space Antenna and provide strong motivation for
developing accurate perturbative descriptions of compact objects moving
in Kerr spacetime~\cite{AmaroSeoane:2012km,Pound:2021qin}.  The Teukolsky
formalism~\cite{Teukolsky:1973ha} is particularly well suited
to this problem: the extreme Weyl scalars are gauge invariant at linear
order and satisfy separable equations, allowing highly accurate
calculations in the frequency and time domains.  Preserving phase accuracy
over an inspiral, however, requires the self-force to post-adiabatic order,
and therefore second-order effects, which are sourced by the metric
perturbation itself rather than by the Weyl scalars~\cite{Pound:2015wva}.
A practical scheme for
metric reconstruction must therefore provide not only a solution of the
linearized Einstein equation, but one with sufficiently good local and
asymptotic gauge properties to construct a well defined second order source.

\subsection{GHZ reconstruction and the asymptotic problem of standard radiation gauges}\label{sec:radiation gauge problem}

The Green--Hollands--Zimmerman (GHZ) reconstruction
formalism~\cite{Green:2019nam} extends
the Chrzanowski--Cohen--Kegeles--Ori
construction~\cite{Chrzanowski:1975wv,Cohen:1974cm,Kegeles:1979an,Ori:2002uv}
to the sourced Einstein equation.  In the IRG-GHZ puncture scheme, the metric perturbation
is decomposed as
\begin{equation}\label{eq:h GHZ}
h^{\rm ret}_{ab} =
	x_{ab}
	+ 2\Re\left( \S^{\dagger}_{ab}\Phi \right)
	+ \dot g_{ab} + \Lie_{\xi}g_{ab},
\end{equation}
Here $\Phi$ is a Hertz potential, satisfying both the $s=-2$
Teukolsky equation and an inversion relation that
ties it to the $s=+2$ Weyl scalar $\psi_{0}$; $x_{ab}$ is a corrector tensor
obeying a triangular system of transport equations along the outgoing
principal null congruence; and the remaining pieces $\dot g_{ab}$ and
$\Lie_{\xi}g_{ab}$ carry the mass, angular-momentum and gauge
perturbations~\cite{Merlin:2016boc,vandeMeent:2017fqk}.
To remove the singular behaviour one further decomposes the retarded field
into a puncture and a residual field~\cite{Toomani:2021jlo},
\begin{equation}
	h^{\rm ret}_{ab}
		=
		h^{P}_{ab}
		+
		h^{R}_{ab},
\end{equation}
where the puncture reproduces the local singular structure of the
retarded field near the worldline.  Multiplying the puncture by a window
function renders its effective source of compact radial support.

The asymptotic behavior of the residual perturbation
is therefore determined solely by the Hertz potential
\beq
h^{R}_{ab}
=
2\Re\left(\S^{\dagger}\Phi^R\right)_{ab},
\qquad r\gg r_\p.
\eeq

This leads to an important practical problem.  In frequency-domain
implementations the Hertz potential is commonly obtained from the
computed Weyl scalar $\psi_0$ using the inversion procedure of
Ori~\cite{Ori:2002uv}.
For an outgoing solution at future null infinity,
\begin{equation}\label{eq:psi0 asy}
\psi_0
=
\frac{\psifive(u,\theta,\varphi)}{r^5}
+O(r^{-6}),
\end{equation}
whereas the corresponding retarded Hertz potential obtained from the
radial inversion
\beq\label{eq:inv rel}
	\th^4 \bar\Phi =2 \psi_0
\eeq
behaves generically as
\begin{equation}
\Phi = O(r^3).
\end{equation}
This growth is familiar from radiation-gauge reconstruction and does
not signal a divergence of the physical curvature.  It does, however,
have an immediate consequence for the reconstructed metric.  Applying
the second-order reconstruction operator $\S^{\dagger}$ to
such a Hertz potential gives, in the Kinnersley Newman--Penrose frame,
the nonvanishing radiation-gauge components schematically as
\begin{equation}
h_{nn},\qquad h_{nm},\qquad h_{mm}
=
O(r),
\end{equation}
with the corresponding complex-conjugate components.  Thus the metric
representative obtained directly from the Ori-inverted Hertz potential
is not asymptotically flat in the Bondi sense\newtext{, and the
quadratic source $S^{(2)}_{ab}[h,h]$ built from it fails to be integrable
against the volume element at $\scri$, so that neither the second-order
Teukolsky source nor the asymptotic flux formulas can be evaluated on it}.

\subsection{Asymptotically flat reconstruction}
\label{sec:af reconstruction}
What is needed instead is a representative of the same physical perturbation
whose Newman--Penrose components decay at the Bondi rate.  Integrability of
the second-order source requires at least
\begin{equation}\label{eq:falloff radiative}
h_{nn},\ h_{nm}, \ h_{mm}
= O(r^{-1}),
\end{equation}
together with
\begin{equation}\label{eq:falloff nonradiative}
h_{ll},\ h_{lm},\ h_{ln},\ h_{m\bar m}
=
O(r^{-2}).
\end{equation}
The discrepancy between the required falloff and what the CCK--Ori
procedure delivers is therefore two powers of $r$ in the radiative
Newman--Penrose components.
\PZ{Worth stating precisely \emph{where}
\eqref{eq:falloff radiative}--\eqref{eq:falloff nonradiative} come from ---
Bondi gauge, or integrability of $S^{(2)}$, or both?  Hollands--Toomani
Sec.~1.3 and Sec.~3.3.2 have the statement you want.}

The $O(r)$ behavior is a gauge effect rather than a failure of the
Teukolsky solution.  Indeed, an asymptotically flat metric may itself
be represented using a Hertz potential with an $O(r^3)$ tail, provided
that the reconstructed piece is supplemented by the appropriate
residual gauge transformation.  Hollands and Toomani~\cite{Hollands:2024iqp} gave a
systematic construction of precisely this transformation.  They showed that the
unwanted asymptotic terms can be removed by combining a residual gauge
vector $\zeta^a$ with a homogeneous modification of the Hertz
potential of the form
\begin{equation}\label{eq:Phi completion intro}
\Phi
\longrightarrow
\Phi
+
\PhiN{0}
+
\frac{\PhiN{1}}{\rho}
+
\frac{\PhiN{2}}{\rho^2}
+
\frac{\PhiN{3}}{\rho^3}.
\end{equation}
The apparently growing terms in this homogeneous addition are essential:
their contribution under $\S^{\dagger}$ combines with
$\Lie_{\zeta}g_{ab}$ to cancel the growing terms in the original
radiation-gauge metric, leaving a perturbation with Bondi-type falloff
at future null infinity~\cite{Hollands:2024iqp}.

In this work we turn this formal asymptotic construction into a
practical numerical procedure suitable for GHZ puncture calculations.
Starting from the asymptotic coefficient $\psifive$ of the
residual Weyl scalar, we determine the Held scalar $\fHo$ that
generates both the homogeneous Hertz completion and the residual gauge
vector.  In Schwarzschild spacetime the construction becomes diagonal
in spin-weighted spherical harmonics, allowing all completion fields to
be obtained explicitly mode by mode.  We implement these relations
numerically and demonstrate the cancellation of the growing
radiation-gauge terms and the recovery of the expected asymptotic
falloff.  We then extend the construction to Kerr, where the input
Teukolsky solution is naturally expressed in spin-weighted spheroidal
harmonics, and present numerical results using an $\emm$-mode
implementation.

The resulting procedure complements the local regularization provided
by the Teukolsky puncture scheme.  The puncture controls the singular
structure in a neighborhood of the compact object, while the
asymptotic completion controls the behavior of the reconstructed field
at null infinity.  Together, these ingredients provide a metric
perturbation with substantially better global properties for
self-force calculations and, in particular, for the construction of
higher-order perturbations in Kerr spacetime.

\subsection{Conventions, notation, and outline}\label{sec:conventions}
% ---------------------------------------------------------------------
% ROADMAP.  This subsection did not exist.  It is worth adding because the
% notebooks carry three separate sign/normalisation conventions that a
% reader will otherwise trip over:
%   (a) signature.  The psi_0 data imported from the Teukolsky package is
%       MOSTLY PLUS (-,+,+,+);  the GHZ/Held expressions, Vahid's notes and
%       the notebooks are all MOSTLY MINUS (+,-,-,-).  The notebook applies
%       flip-psi0Convention: psi_0 -> -psi_0 and Phi -> -Phi on import.
%   (b) the "Re" of a mode:  Re[X]_{lm} = 1/2 (X_{lm} + (-1)^{m+s} Xbar_{l,-m})
%   (c) the factor 2 in h = 2 Re(S^dagger Phi)
% Write these out once, here, and refer back.
% ---------------------------------------------------------------------
We work in geometrized units $G=c=1$ and use the metric signature
$(+,-,-,-)$.  Newman--Penrose and GHP conventions follow
GHZ~\cite{Green:2019nam}, and Held's modified GHP
calculus~\cite{Held:1974,Held:1975} follows
Ref.~\cite{Hollands:2024iqp}, with
$\Held{\ \cdot\ }$ denoting Held (``$r$-independent'') scalars.
The Kinnersley tetrad is used throughout, so that
\begin{equation}\label{eq:rho def}
\rho = -\frac{1}{r-ia\cos\theta},
\qquad
\bar\rho = -\frac{1}{r+ia\cos\theta},
\end{equation}
and the nonvanishing background Held scalars are
\begin{equation}\label{eq:bkgd held}
\PsiH = M,\quad
\rhopH = -\tfrac12,\quad
\tauH = - ia\sin\theta/\sqrt2,\quad
\Oh = -2ia\cos\theta ,
\end{equation}
together with their conjugates.  Note that $\Oh$ is a spin-weight-zero multiplication operator
proportional to $\cos\theta$, so it is represented in the
$\sYY{s}{\l}{\emm}$ basis by the same tridiagonal matrix $\mathsf C$ that
represents $\rho^{-1}$ in \eqref{eq:R matrix}; explicitly
$\Oh = -2ia\,\mathsf C$.


Four conventions recur and are fixed once here.
\begin{enumerate}[label=(\roman*)]
\item \emph{Signature of the imported Weyl data.}  All formulas in this
paper are written in the mostly-minus signature $(+,-,-,-)$, as are the GHZ
and Held expressions they are built from.  Teukolsky amplitudes produced with
the Black Hole Perturbation Toolkit~\cite{BHPToolkit}, however, are in the
mostly-plus convention $(-,+,+,+)$, and the two are related by
\begin{equation}\label{eq:signature flip}
\psi_0^{(+---)}=-\psi_0^{(-+++)},
\qquad
\Phi^{(+---)}=-\Phi^{(-+++)} .
\end{equation}
The imported amplitudes are therefore negated once, on import; this is the
origin of the minus sign in Step~1 of Sec.~\ref{sec:algorithm}.

\item \emph{Conjugation of a spin-weighted field.}  Spin-weighted spherical
harmonics obey
\begin{equation}\label{eq:sYlm conj}
\overline{\sYY{s}{\l}{\emm}(\theta,\varphi)}
=
(-1)^{s+\emm}\;\sYY{-s}{\l}{-\emm}(\theta,\varphi),
\end{equation}
so that if $X$ has spin weight $s$ and coefficients $X_{\l\emm}$ in the
$\sYY{s}{\l}{\emm}$ basis, then $\bar X$ has spin weight $-s$ and its
coefficients \emph{in the $\sYY{-s}{\l}{\emm}$ basis} are
\begin{equation}\label{eq:conj mode}
\big(\bar X\big)_{\l\emm}
=
(-1)^{s+\emm}\;\overline{X_{\l,-\emm}}\,,
\end{equation}
where the overline on the right is ordinary complex conjugation of the mode
amplitude, accompanied by $\omega\to-\omega$ (equivalently, by
$(\emm,n,k)\to(-\emm,-n,-k)$, since $\wmnk$ is odd under that reflection).
The factor $(-1)^{s}$ is easy to lose because it is unity for every
even-spin object in this construction --- $\psifive$, $\psifourone$, $\fHo$
and all four $\PhiN{n}$ have $s=\pm2$ --- but it equals $-1$ for the
odd-spin ones: $\zetam$, $\zetamb$, the intermediate chains
$\edthH\fbHo$ and $\edthpH\edthH^{2}\fbHo$, and the metric components
$h_{nm}$, $h_{n\bar m}$.  See Table~\ref{tab:spin weights}.

\item \emph{Real part.}  A real part is meaningful in mode space only for
$s=0$, since $X$ and $\bar X$ otherwise live in different harmonic bases.
For a real spin-zero scalar, \eqref{eq:conj mode} gives
\begin{equation}\label{eq:Re mode}
\big(\Re X\big)_{\l\emm}
=
\tfrac12\Big(X_{\l\emm} + (-1)^{\emm}\,\overline{X_{\l,-\emm}}\Big),
\end{equation}
which is the projection used for $\zetal$ and $\zetan$ in
\eqref{eq:zeta Held}.  For a component pair of nonzero spin weight, such as
$h_{nm}$ and $h_{n\bar m}$, reality of $h_{ab}$ is instead the statement
$h_{nm}=\overline{h_{n\bar m}}$, which in modes is \eqref{eq:conj mode} with
$s=-1$ and therefore carries an explicit minus sign.

\item \emph{Reconstruction normalization.}  We write the reconstructed
piece as $h_{ab}=2\Re(\S^\dagger\Phi)_{ab}$, so that the operator
$\S^{\dagger}$ quoted in Appendix~\ref{app:Sdag} is \emph{not} already
real-projected.
\PZ{This is the one place where the notebook and Vahid's notes differ by a
uniform factor of 2 in all four $\PhiN{n}$.  Pick one and state it here;
the equations below use the notes' normalisation.}
\end{enumerate}

The paper is organized as follows.  Section~\ref{sec:prelim} collects the
GHZ decomposition and the Held formalism in the form we need.
Section~\ref{sec:adjusted} contains the main analytic result: the adjusted
reconstruction, presented as an algorithm generated by a single Held seed
$\fHo$.  Section~\ref{sec:modes} reduces every step of that algorithm to
linear algebra on spin-weighted spherical-harmonic coefficients.
Sections~\ref{sec:schw} and \ref{sec:kerr} present the Schwarzschild and Kerr
implementations and numerical results, and Section~\ref{sec:discussion}
concludes.  Appendices collect the Held operator matrix elements, the
Teukolsky--Starobinsky constants, and explicit mode expressions for
$\S^{\dagger}$ and $\Lie_\zeta g$.


% =====================================================================
\section{Preliminaries}
\label{sec:prelim}
% =====================================================================
% ---------------------------------------------------------------------
% ROADMAP.  Keep this section short and purely expository -- it exists so
% the paper is self-contained for a self-force reader who has not read
% Hollands-Toomani.  Three subsections, ~1 page each at most.
% ---------------------------------------------------------------------

\subsection{Kerr background, NP frame, and Held scalars}
\label{sec:held}
% ROADMAP: Kinnersley tetrad; GHP weights; Held's "integration-constant"
% calculus: given a GHP scalar eta of boost/spin weight (p,q), Held's
% tilded operators commute with the radial integration, so a general
% solution of a thorn-equation is written as a Held scalar times powers of
% rho.  State the dictionary
%       thorn-tilde' = ... ,  edth-tilde = ... ,  and on a Fourier mode
%       thorn-tilde' -> -i omega.
% That last replacement is what makes the whole scheme algebraic and is
% worth displaying.
Held's modification of the GHP formalism~\cite{Held:1974,Held:1975} is what
makes the completion algebraic rather than differential, and we recall only
as much of it as we need.  In a type-D background one may trade the GHP
operators for tilded counterparts $\thpH$, $\edthH$, $\edthpH$ that commute
with radial integration along the outgoing principal null congruence.  The
general solution of a $\th$-equation is then a sum of powers of $\rho$ with
coefficients that are annihilated by $\th$ --- the \emph{Held scalars},
written $\Held{\ \cdot\ }$ --- so that the integration constants of the
radial problem become fields on the sphere.  The four quantities $\PhiN{n}$
of \eqref{eq:Phi completion intro} are exactly such integration constants,
which is why the completion is fixed by angular data alone.

On a fixed Fourier mode $e^{-i\omega u}$ at $\scri$ the Held operator
$\thpH$ acts algebraically,
\begin{equation}\label{eq:thorn prime mode}
\thpH \longrightarrow -i\omega ,
\end{equation}
which is the single fact that converts the transport equations of
Sec.~\ref{sec:adjusted} into algebraic relations.
\PZ{State the precise sense of \eqref{eq:thorn prime mode}: $\thpH$ is
$\pd_u$ acting on the Held (i.e.\ $r$-independent) part at $\scri$.}

The spin weights of every object appearing in the construction are collected
in Table~\ref{tab:spin weights}.  They fix both the angular matrix elements
of Sec.~\ref{sec:held operators} and, through \eqref{eq:conj mode}, the sign
attached to each complex conjugation; the odd-spin entries are the ones for
which the factor $(-1)^{s}$ is not unity.

\begin{table}[htb]
\begin{ruledtabular}
\begin{tabular}{lcccccccc}
 & $\fHo$ & $\fbHo$ & $\PhiN{n}$ & $\edthH\fbHo$ & $\edthH^{2}\fbHo$
 & $\zetal,\ \zetan$ & $\zetam$ & $\zetamb$ \\
\hline
spin weight $s$ & $+2$ & $-2$ & $-2$ & $-1$ & $0$ & $0$ & $+1$ & $-1$\\
$(-1)^{s}$ & $+$ & $+$ & $+$ & $-$ & $+$ & $+$ & $-$ & $-$\\
\end{tabular}
\end{ruledtabular}
\begin{tabular}{lccccccc}
 & $h_{ll}$ & $h_{ln}$ & $h_{nn}$ & $h_{m\bar m}$ & $h_{nm}$ & $h_{n\bar m}$
 & $h_{mm}$\\
\hline
spin weight $s$ & $0$ & $0$ & $0$ & $0$ & $+1$ & $-1$ & $+2$\\
$(-1)^{s}$ & $+$ & $+$ & $+$ & $+$ & $-$ & $-$ & $+$\\
\end{tabular}
\caption{Spin weights of the fields entering the adjusted reconstruction.
The second row of each block is the sign $(-1)^{s}$ appearing in the
conjugation rule \eqref{eq:conj mode}; it is $-1$ precisely for the
odd-spin-weight quantities.}
\label{tab:spin weights}
\end{table}
\PZ{The field spin weights are the values hard-coded as
\texttt{Spinof[...]} in the notebooks.  The metric-component spin weights
follow from the GHP type of $m^{a}$ and $n^{a}$; please double-check the
sign convention for $h_{nm}$ versus $h_{n\bar m}$ against your own, since the
two appear swapped between \texttt{Re$\S$dagMode$\Phi$Ops} and
\texttt{LieZetagComps} (see Appendices~\ref{app:Sdag} and \ref{app:Lie}).}

\subsection{The GHZ decomposition and the puncture scheme}
\label{sec:ghz}
% ROADMAP: restate (1.1) properly, say what each piece solves, and state
% the puncture split and the window function.  Then make the key point in
% one sentence: because the effective source is compactly supported in r,
% the field near scri is *purely* 2 Re(S^dag Phi^R), so the asymptotic
% problem is entirely a statement about Phi.
Following Ori, take a slice $\Sigma$ of constant Boyer--Lindquist time
$t=t_{0}$ and split the spacetime into the regions to the left and right of
the sphere $\mathcal S_{\p}$ at $r=\rp$, denoted $\Sigma_{-}$ and
$\Sigma_{+}$ respectively.  In the world-tube region
$\Gamma = \{\rmn\leq r \leq \rmx\}$ define $\Gamma_{\pm} = \Gamma \cap
\Sigma_{\pm}$, and let the vacuum regions $\Sigma_{\pm}'$ be the complements
of $\Gamma_{\pm}$ in $\Sigma_{\pm}$.  Define $\Phi^{\pm}$ to be the regular
(vacuum) solutions supported on $\Sigma_{\pm}'$, so that
\begin{equation}
h^{\pm}_{ab} = 2 \Re ( \S_{ab}^{\dagger} \Phi^{\pm}), \qquad x \in \Sigma_{\pm}' .
\end{equation}
The asymptotic decay of $h^{+}_{ab}$, and therefore of the residual field, is
completely determined by $\Phi^{+}$, which is in turn determined by the decay
of $\psi_{0}$.

\subsection{Asymptotics at $\scri$ and the required falloff}
\label{sec:asymptotics}
% ROADMAP: make explicit that you are *not* imposing full Bondi gauge, only
% the falloff conditions needed for the second-order source.  A short remark
% on what is left over (supertranslation freedom) belongs here or in the
% Discussion.
The target of the construction is the pair of conditions
\eqref{eq:falloff radiative}--\eqref{eq:falloff nonradiative}.  Two
observations make them easier to check.  First, in the ingoing radiation
gauge the non-radiative components $h_{ll}$, $h_{ln}$, $h_{lm}$,
$h_{l\bar m}$ and $h_{m\bar m}$ vanish \emph{identically}, both for the
reconstructed piece (Appendix~\ref{app:Sdag}) and for the gauge term
(Appendix~\ref{app:Lie}); condition \eqref{eq:falloff nonradiative} is
therefore satisfied trivially, and the whole problem lives in the five
radiative components.  Second, because the Hertz potential enters through
$\S^{\dagger}$, which is second order in radial derivatives, each radiative
component is a finite Laurent polynomial in $r$ near $\scri$; verifying
\eqref{eq:falloff radiative} reduces to showing that two coefficients ---
those of $r$ and of $r^{0}$ --- vanish in each of them.  That is the check
carried out in Sec.~\ref{sec:schw verification}.

\PZ{Decide whether to call the end result ``Bondi gauge'' or
``asymptotically flat gauge''.  Strictly you fix the falloff, not the gauge;
the residual freedom is the usual BMS supertranslation plus the
$\ell=0,1$ modes.  Hollands--Toomani Sec.~3.3.2 relates $\PhiN{n}$ to the
news and memory tensors, which would make a nice remark here.}


% =====================================================================
\section{Adjusted asymptotically flat GHZ reconstruction}
\label{sec:adjusted}
% =====================================================================
% ---------------------------------------------------------------------
% ROADMAP.  This is the analytic core.  The narrative to aim for:
%   "Everything -- all four Hertz integration constants and all four
%    components of the residual gauge vector -- is generated by ONE Held
%    scalar f^o, and f^o is obtained from psi_0^{5o} by inverting a single
%    eighth-order angular operator, which in modes is just division by the
%    Teukolsky-Starobinsky constant."
% That sentence is the paper.  Sections 3.1-3.5 unpack it, 3.6 boxes it.
% ---------------------------------------------------------------------

\subsection{Asymptotic Weyl amplitudes}
\label{sec:weyl amplitudes}
% ---------------------------------------------------------------------
% NOTATION.  One symbol for the Teukolsky boundary amplitudes throughout:
%   {}_s Z^infty  (at scri)  and  {}_s Z^H  (at the past horizon).
% Spin weight in the left subscript, boundary in the superscript.  This
% replaces the notebooks' C^{+}/C^{-}, where the +/- denoted the boundary
% and the spin was carried by a separate argument -- a clash with the
% +/- of {}_{\pm2}Z that is easy to trip over.  The dictionary is
%   C^{+}_{lmnk} (s = +2)  ==  {}_{+2}Z^infty_{lmnk}
%   C^{-}_{lmnk} (s = +2)  ==  {}_{+2}Z^{H}_{lmnk}
% ---------------------------------------------------------------------
We write $\Zinf{s}{\lmnk}$ and $\Zhor{s}{\lmnk}$ for
the asymptotic amplitudes of the spin-$s$ Teukolsky solution at future null
infinity and at the past horizon respectively; these are precisely the
$\mathscr I$ and $\mathrm H$ amplitudes returned by any frequency-domain
Teukolsky solver, so that no new radial integrations are needed anywhere in
what follows.  Only the $\scri$ amplitudes enter the asymptotic completion.
The asymptotic coefficients
$\psifiveM{\lmnk}\defeq\Zinf{+2}{\lmnk}$ in
\beq\label{eq:psi0 5o expansion}
\psifive(u,\theta,\varphi) \defeq
 	\sum_{\lmnk}
	  \psifiveM{\lmnk}\,
	 \sSS{+2}{\l}{\emm}(\theta,\varphi;a\wmnk)\,
	 e^{-i\wmnk u}
\eeq
and $\psifouroneM{\lmnk}\defeq\Zinf{-2}{\lmnk}$ in
\beq\label{eq:psi4 1o expansion}
\psifourone(u,\theta,\varphi)  \defeq
	\sum_{\lmnk}
	 \psifouroneM{\lmnk}\,
	 \sSS{-2}{\l}{\emm}(\theta,\varphi;a\wmnk)\,
	  e^{-i\wmnk u}
\eeq
are not independent.  The Teukolsky--Starobinsky identities
\beq \label{eq:TS NP form}
\Delta^{2} \mathcal D^4 (\Delta^{2}\psi_{0,\lmnk}) =
	\frac{1}{4}\mathcal L^{4}(\zeta^{4}\psi_{4,\lmnk})
	+ 3 M \pd_{t}(\zeta^{4}\bar\psi_{4,\lmnk})
\eeq
together with the angular inversion relation
\beq\label{eq:L4 on S}
\mathcal L^{4} \sSS{-2}{\l}{\emm} = D_{\l\emm} \, \sSS{+2}{\l}{\emm}
\eeq
imply
\beq\label{eq:Zplus from Zminus}
\Zinf{+2}{\lmnk} =
	\frac{1}{4 \omega^{4}}
	\left(
	D_{\l\emm}-12iM\omega\, p_{\lmnk}
	\right)
	\Zinf{-2}{\lmnk},
\eeq
where $p_{\lmnk}=\pm1$ is a parity factor related to the phase conventions of
the spin-weighted spheroidal harmonics.
Equivalently, in the form used by the implementation,
\begin{equation}\label{eq:psi4 from psi0}
\psifouroneM{\l\emm}(\omega)
=
\frac{\omega^{4}}{\Dlw}
\Big[
\BcoK{-2}{\l,-\emm}{-a\omega}^{\!*}\,\psifiveM{\l\emm}(\omega)
+3i\omega(-1)^{\emm}\,\bar\psi^{5\circ}_{0,\l,-\emm}(-\omega)
\Big],
\end{equation}
with
\begin{equation}\label{eq:Dlw def}
\Dlw
\defeq
\BcoK{-2}{\l\emm}{a\omega}\,\overline{\BcoK{-2}{\l,-\emm}{-a\omega}}
+9M^{2}\omega^{2}.
\end{equation}
% ROADMAP: this is the paragraph that sells the paper to a practitioner.
% Spell out that {}_{-2}Z^infty is exactly the quantity every frequency-domain
% Teukolsky code already returns (the "I" amplitude in the Toolkit), so the
% completion costs nothing beyond post-processing.
Here $D_{\l\emm}$ is the real angular Teukolsky--Starobinsky constant of
Appendix~\ref{app:TS}, and $\BcoK{-2}{\l\emm}{a\omega}$ is the eigenvalue
combination in which the completion formulas are most compactly written.
\PZ{These two are the same object up to a factor: in Schwarzschild
$\Bco{2}=\LamL/4$ and $D_{\l}=\LamL$, so $\Bco{2}=D_{\l}/4$ and hence
$\Cts\bar\Cts = D^{2}+144M^{2}\omega^{2} = 16\,\Dlw$.  Please confirm that
the same identification holds in Kerr; if it does, one of the two symbols can
be eliminated from the paper entirely, as we did for $C^{\pm}$ versus $Z$.}

For a generic bound geodesic, $\wmnk = \emm\Omega_{\varphi}+n\Omega_{r}+k\Omega_{\theta}$
and the polar reflection symmetry of the spheroidal functions gives
\begin{equation}\label{eq:Z symmetry}
\Zinf{s}{\l,-\emm,-n,-k} = (-1)^{\l+k}\,\big(\Zinf{s}{\lmnk}\big)^{*},
\end{equation}
which halves the number of modes that must be computed.
\PZ{The notebooks apply \eqref{eq:Z symmetry} to $s=+2$ only; state whether
it is meant to hold for general $s$ (it should, being a property of the
spheroidal functions under $\theta\to\pi-\theta$) and cite the source.}

\subsection{The Held seed $\fHo$}
\label{sec:seed}
% ROADMAP: introduce f^o as *the* generator.  One paragraph: it is the Held
% scalar whose eighth-order angular image is the asymptotic Weyl amplitude.
% Emphasise that (3.x) is an ODE in u only through thorn', so on a Fourier
% mode it is purely angular -> purely algebraic in modes.
The entire completion is generated by a single Held scalar $\fHo$ of spin
weight $+2$, defined implicitly by
\begin{equation}\label{eq:f seed}
2\,\thpH{}^{3}\,\bar\psi^{5\circ}_{0}
=
\edthpH{}^{4}\,\edthH^{4}\,\fbHo
-9\,\Psih\Psibh\,\thpH{}^{2}\,\fbHo .
\end{equation}
On a Fourier mode, using \eqref{eq:thorn prime mode}, Eq.~\eqref{eq:f seed}
becomes the purely angular relation
\begin{equation}\label{eq:f seed mode op}
\edthpH{}^{4}\,\edthH^{4}\,\fbHo + 9M^{2}\omega^{2}\,\fbHo
=
2i\omega^{3}\,\bar\psi^{5\circ}_{0},
\end{equation}
which is diagonalized by the spin-weighted spheroidal harmonics.  Its
solution is
\begin{equation}\label{eq:f mode kerr}
\boxed{\;
\fHo_{\lmnk}
=
\frac{2i\,\wmnk^{3}}{\Dlw}\;\psifiveM{\lmnk}
\;=\;
\frac{32\,i\,\wmnk^{3}}{\Cts_{\l\emm}\overline{\Cts_{\l\emm}}}\;\psifiveM{\lmnk}
\;}
\end{equation}
where $\Cts_{\l\emm}$ is the complex Teukolsky--Starobinsky constant of
Appendix~\ref{app:TS}.
\PZ{Check which of the two right-hand sides of \eqref{eq:f mode kerr} the
notebook's \texttt{Starobinsky$\mathcal C$} actually returns: the
\texttt{CTeukolskyStarobinsky} routine returns the \emph{complex} $\Cts=D+12iP\tilde\omega$,
but \texttt{fHS} divides $32i\omega^3\psifive$ by it directly, which only
reproduces the Schwarzschild limit $2i\omega^3/(\Bco{2}^2+9M^2\omega^2)$ if
the symbol actually holds $\Cts\bar\Cts$.  One of the two is off; please
confirm in the notebook.}

\subsection{Held integration constants $\PhiN{n}$}
\label{sec:Phi n}
% ROADMAP: present (3.x)-(3.y) as a *triangular* system -- Phi^{3o} is
% algebraic in f, Phi^{2o} is algebraic, Phi^{1o} needs one thorn'-inversion,
% Phi^{0o} needs two.  On a Fourier mode each inversion is division by
% -i omega, so nothing is ever integrated numerically.  Say that explicitly;
% it is the practical punchline of this subsection.
The four Held integration constants appearing in
\eqref{eq:Phi completion intro} are determined from $\fbHo$ by the triangular
system
\begin{subequations}\label{eq:Phi n Held}
\begin{align}
\PhiN{3} &= -2\,\thpH\fbHo,
\label{eq:Phi3 Held}\\[2pt]
\PhiN{2} &= -2\,\edthpH\edthH\fbHo,
\label{eq:Phi2 Held}\\[2pt]
\thpH\,\PhiN{1} &= -\edthpH{}^{2}\edthH^{2}\fbHo
	-\big(4\tabh\,\edthH+3\Psih\big)\thpH\fbHo,
\label{eq:Phi1 Held}\\[2pt]
\thpH{}^{2}\,\PhiN{0} &= -\tfrac13\,\edthpH{}^{3}\edthH^{3}\fbHo
	-\big[2\tabh\,\edthpH\edthH^{2}
	 +\Psih\,\edthpH\edthH
	 +2\Psih\rhopH\big]\thpH\fbHo .
\label{eq:Phi0 Held}
\end{align}
\end{subequations}
On a Fourier mode $\thpH\to-i\omega$ and the two inversions in
\eqref{eq:Phi1 Held}--\eqref{eq:Phi0 Held} are divisions by $-i\omega$ and
$-\omega^{2}$ respectively; no radial or temporal integration is required.
The system is therefore purely angular, and reduces in modes to the banded
linear algebra of Sec.~\ref{sec:modes}.
\PZ{$\PhiN{0}$ and $\PhiN{3}$ are called $\dHo$ and $\eHo$ in Vahid's notes
(Eqs.~29 and 32); I have used $\PhiN{n}$ throughout.}

\subsection{The residual gauge vector $\zeta^{a}$}
\label{sec:zeta}
% ROADMAP: same structure: zeta_mbar is algebraic in f, zeta_l and zeta_n
% each require one thorn'-inversion.  Note the Re{} projection: zeta_l and
% zeta_n are real spin-0 scalars, so in modes they use (ii) of Sec 1.4.
The residual gauge vector is likewise generated by $\fbHo$:
\begin{subequations}\label{eq:zeta Held}
\begin{align}
\zetamb &= \edthH\fbHo,
\label{eq:zeta m Held}\\[2pt]
\thpH\,\zetal &= -\Re\big\{\edthH\edthH\fbHo\big\},
\label{eq:zeta l Held}\\[2pt]
\thpH\,\zetan &= -\Re\Big\{
	\edthpH\edthH^{3}\fbHo
	-\rhopH\,\edthH\edthH\fbHo
	-2\tah\,\edthH\thpH\fbHo
	-\tfrac12\Oh\,\edthH\edthH\thpH\fbHo
	\Big\}.
\label{eq:zeta n Held}
\end{align}
\end{subequations}
Here $\zetal$ and $\zetan$ are real scalars of spin weight zero, so that
$\Re$ is implemented in modes through \eqref{eq:Re mode}.  The conjugate of
\eqref{eq:zeta m Held} is best written directly as $\zetam=\edthpH\fHo$
rather than obtained by conjugating $\zetamb$; see Sec.~\ref{sec:reality}.

\subsection{The adjusted metric and cancellation of the growing terms}
\label{sec:adjusted metric}
% ROADMAP: state the theorem-level result: with (3.x) and (3.y),
%    h^adj = 2 Re(S^dag Phi^adj) - L_zeta g
% satisfies the falloff (1.8)-(1.9).  Then say how you *verify* it: expand
% each NP component in powers of r, show the O(r) and O(r^0) coefficients
% vanish identically mode by mode.  That is exactly what hAdjustment /
% Coefficient[...,r] does in the notebooks, and it is a real result worth
% a short table.
Combining the completed Hertz potential
\begin{equation}\label{eq:Phi adj}
\Phiadj = \Phi + \PhiN{0} + \frac{\PhiN{1}}{\rho} + \frac{\PhiN{2}}{\rho^{2}}
+ \frac{\PhiN{3}}{\rho^{3}}
\end{equation}
with the residual gauge transformation gives the adjusted metric
perturbation
\begin{equation}\label{eq:h adj}
\boxed{\;
h^{\rm adj}_{ab} = 2\Re\big(\S^{\dagger}_{ab}\Phiadj\big) - \Lie_{\zeta}g_{ab}.
\;}
\end{equation}
The Lie-derivative term has the same ingoing-radiation-gauge structure as the
reconstructed piece,
\begin{equation}\label{eq:Lie zeta IRG}
\Lie_{\zeta}g_{ab}\,l^{a} = 0,
\qquad
\Lie_{\zeta}g_{ab}\,m^{a}\bar m^{b} = 0,
\end{equation}
so that the adjustment preserves the IRG conditions; its nonvanishing
components are collected in Appendix~\ref{app:Lie}.

That \eqref{eq:h adj} is asymptotically flat is a theorem of
Ref.~\cite{Hollands:2024iqp}.  What the mode decomposition adds is an
explicit, checkable form of it: with $\fHo$ given by \eqref{eq:f mode kerr},
$\PhiN{n}$ by \eqref{eq:Phi n Held} and $\zeta^{a}$ by
\eqref{eq:zeta Held}, the coefficients of $r$ and of $r^{0}$ in each of
$h^{\rm adj}_{nn}$, $h^{\rm adj}_{nm}$, $h^{\rm adj}_{n\bar m}$,
$h^{\rm adj}_{mm}$ and $h^{\rm adj}_{\bar m\bar m}$ cancel mode by mode
between the two terms of \eqref{eq:h adj}, leaving the Bondi falloff
\eqref{eq:falloff radiative}; the non-radiative components
\eqref{eq:falloff nonradiative} vanish identically in the ingoing radiation
gauge, as noted in Sec.~\ref{sec:asymptotics}.  We verify this cancellation
symbolically in Sec.~\ref{sec:schw verification}.

\subsection{Summary of the algorithm}
\label{sec:algorithm}
% ROADMAP: the boxed recipe.  Keep it to one page and make it copy-able:
% a reader should be able to implement the paper from this list alone.
\begin{enumerate}[label=\textbf{Step \arabic*.},leftmargin=*]
\item Solve the $s=+2$ Teukolsky equation with the point-particle source and
extract the asymptotic amplitudes $\Zinf{+2}{\lmnk}$ at $\scri$; set
$\psifiveM{\lmnk} = -\Zinf{+2}{\lmnk}$ (signature flip, Sec.~\ref{sec:conventions}).
\item Form the seed $\fHo_{\lmnk}$ from \eqref{eq:f mode kerr}: one complex
division per mode.
\item Project $\fHo$ from the spheroidal $\sSS{+2}{\l}{\emm}$ basis onto the
spherical $\sYY{+2}{L}{\emm}$ basis using the mixing coefficients
\eqref{eq:spheroidal to spherical}.  (In Schwarzschild this step is the
identity.)
\item Build $\PhiN{0},\dots,\PhiN{3}$ from \eqref{eq:Phi n Held} by repeated
application of the tridiagonal operators of Sec.~\ref{sec:held operators}.
\item Assemble $\Phiadj$ from \eqref{eq:Phi adj}; the inverse powers of
$\rho$ are the tridiagonal $\cos\theta$ matrix \eqref{eq:R matrix}, so this
step widens the $\l$-support by three and requires three extra $\l$'s of
padding.
\item Build $\zetal,\zetan,\zetam$ from \eqref{eq:zeta Held} and evaluate
$\Lie_{\zeta}g_{ab}$ from Appendix~\ref{app:Lie}.
\item Assemble $h^{\rm adj}_{ab}$ from \eqref{eq:h adj} and verify the
falloff.
\end{enumerate}
\PZ{A figure here -- a simple flow chart from $\Zinf{+2}{}$ to $h^{\rm adj}$ --
would carry a lot of weight for a methods paper.  Cheap to make and the
referee will like it.}


% =====================================================================
\section{Mode decomposition}
\label{sec:modes}
% =====================================================================
% ---------------------------------------------------------------------
% ROADMAP.  This section is the actual new technical content and should be
% written carefully; everything in Sec. 3 is a restatement of
% Hollands-Toomani, but Sec. 4 is yours.  Structure:
%   4.1 frequency conventions
%   4.2 spheroidal -> spherical
%   4.3 Held angular operators as tridiagonal matrices   <-- the key result
%   4.4 multiplication by 1/rho -> cos theta matrix       <-- the second one
%   4.5 conjugation / reality
% The take-away sentence: "every operator in Sec. 3 is tridiagonal in l at
% fixed m, so the whole completion is O(l_max) work per (m,n,k)."
% ---------------------------------------------------------------------

\subsection{Frequency-domain conventions}
\label{sec:freq}
For a generic bound geodesic with radial, polar and azimuthal frequencies
$(\Omega_{r},\Omega_{\theta},\Omega_{\varphi})$ we use
\begin{equation}\label{eq:omega mnk}
\wmnk = \emm\,\Omega_{\varphi}+n\,\Omega_{r}+k\,\Omega_{\theta},
\end{equation}
reducing for circular equatorial orbits to $\omega = \emm\Omega_{\varphi}$
with $\Omega_{\varphi}=\sqrt{M/r_{0}^{3}}$ in Schwarzschild.
\PZ{Note that $\fHo\propto\omega^{3}$, so the $\omega=0$ modes drop out of the
radiative completion entirely; the notebook returns an empty association for
them.  The stationary ($\omega=0$) completion -- mass and angular momentum --
is handled by $\dot g_{ab}$ in \eqref{eq:h GHZ}, i.e.\ by the completion of
Refs.~\cite{Merlin:2016boc,vandeMeent:2017fqk}; say so here.}

\subsection{Spheroidal-to-spherical projection}
\label{sec:spheroidal}
The Teukolsky solution is naturally expressed in spin-weighted
\emph{spheroidal} harmonics, whereas the Held operators of
Sec.~\ref{sec:held operators} act simply on spin-weighted \emph{spherical}
harmonics.  We therefore expand
\begin{equation}\label{eq:spheroidal to spherical}
\sSS{s}{\l}{\emm}(\theta,\varphi;a\omega)
=
\sum_{L}\,b^{(s,\l,\emm)}_{L}(a\omega)\;\sYY{s}{L}{\emm}(\theta,\varphi),
\end{equation}
so that a field with spheroidal amplitudes $X_{\l\emm}$ has spherical
amplitudes
\begin{equation}\label{eq:S to Y amplitudes}
X^{(Y)}_{L\emm} = \sum_{\l} X_{\l\emm}\, b^{(s,\l,\emm)}_{L}(a\omega).
\end{equation}
\PZ{State how many terms you keep ($\texttt{NumTerms}\to30$ in the notebook)
and the chop tolerance ($10^{-25}$), and give a one-line convergence
statement -- the coefficients fall off geometrically in $|L-\l|$ for
$a\omega\lesssim1$.  A small table or plot of $|b_L|$ vs $L-\l$ for the
worst-case $a\omega$ in your runs would settle the truncation question
for a referee.}

\subsection{Held angular operators in the $\sYY{s}{\l}{\emm}$ basis}
\label{sec:held operators}
% ---------------------------------------------------------------------
% ROADMAP.  Present these as the central technical tool.  Say in words
% what they mean: in Schwarzschild edth-tilde is diagonal (the usual
% spin-raising eigenvalue); the entire a-dependence of the Held calculus
% enters as a *tridiagonal* correction with a single small parameter
% a omega.  Then note the structural relation edth-tilde = edth-tilde|_{a=0}
% + thorn-tilde' tau-tilde^o, which is a nontrivial check.
% ---------------------------------------------------------------------
Write $f=\sum_{\l\emm}f_{\l\emm}\,\sYY{s}{\l}{\emm}(\theta,\varphi)$ for a
Held scalar of spin weight $s$.  The Held angular operators are then
tridiagonal in $\l$ at fixed $\emm$:
\begin{subequations}\label{eq:held matrices}
\begin{align}
\big(\edthH f\big)_{\l\emm}
&=
\Big(1-\frac{a\omega\,\emm}{\l(\l+1)}\Big)
\sqrt{\tfrac{(\l-s)(\l+s+1)}{2}}\;f_{\l\emm}
\non
&\quad
+\frac{a\omega}{\l+1}
\sqrt{\tfrac{(\l-\emm+1)(\l+\emm+1)(\l-s)(\l-s+1)}{2(2\l+1)(2\l+3)}}\;f_{\l+1,\emm}
\non
&\quad
-\frac{a\omega}{\l}
\sqrt{\tfrac{(\l-\emm)(\l+\emm)(\l+s)(\l+s+1)}{2(2\l-1)(2\l+1)}}\;f_{\l-1,\emm},
\label{eq:edth matrix}\\[6pt]
\big(\edthpH f\big)_{\l\emm}
&=
-\Big(1-\frac{a\omega\,\emm}{\l(\l+1)}\Big)
\sqrt{\tfrac{(\l+s)(\l-s+1)}{2}}\;f_{\l\emm}
\non
&\quad
+\frac{a\omega}{\l+1}
\sqrt{\tfrac{(\l-\emm+1)(\l+\emm+1)(\l+s)(\l+s+1)}{2(2\l+1)(2\l+3)}}\;f_{\l+1,\emm}
\non
&\quad
-\frac{a\omega}{\l}
\sqrt{\tfrac{(\l-\emm)(\l+\emm)(\l-s)(\l-s+1)}{2(2\l-1)(2\l+1)}}\;f_{\l-1,\emm},
\label{eq:edthp matrix}
\end{align}
\end{subequations}
while the multiplication operators by the background Held scalars are
\begin{subequations}\label{eq:mult matrices}
\begin{align}
\big(\tah f\big)_{\l\emm}
&=
-\frac{ia\,\emm}{\l(\l+1)}\sqrt{\tfrac{(\l-s)(\l+s+1)}{2}}\;f_{\l\emm}
\non
&\quad
+\frac{ia}{\l+1}
\sqrt{\tfrac{(\l-\emm+1)(\l+\emm+1)(\l-s)(\l-s+1)}{2(2\l+1)(2\l+3)}}\;f_{\l+1,\emm}
\non
&\quad
-\frac{ia}{\l}
\sqrt{\tfrac{(\l-\emm)(\l+\emm)(\l+s)(\l+s+1)}{2(2\l-1)(2\l+1)}}\;f_{\l-1,\emm},
\label{eq:tau matrix}\\[6pt]
\big(\tabh f\big)_{\l\emm}
&=
+\frac{ia\,\emm}{\l(\l+1)}\sqrt{\tfrac{(\l+s)(\l-s+1)}{2}}\;f_{\l\emm}
\non
&\quad
+\frac{ia}{\l+1}
\sqrt{\tfrac{(\l-\emm+1)(\l+\emm+1)(\l+s)(\l+s+1)}{2(2\l+1)(2\l+3)}}\;f_{\l+1,\emm}
\non
&\quad
-\frac{ia}{\l}
\sqrt{\tfrac{(\l-\emm)(\l+\emm)(\l-s)(\l-s+1)}{2(2\l-1)(2\l+1)}}\;f_{\l-1,\emm},
\label{eq:taubar matrix}\\[6pt]
\big(\Oh f\big)_{\l\emm}
&=
+\frac{2ia\,\emm s}{\l(\l+1)}\;f_{\l\emm}
-\frac{2ia}{\l+1}
\sqrt{\tfrac{(\l-\emm+1)(\l+\emm+1)(\l-s+1)(\l+s+1)}{(2\l+1)(2\l+3)}}\;f_{\l+1,\emm}
\non
&\quad
-\frac{2ia}{\l}
\sqrt{\tfrac{(\l-\emm)(\l+\emm)(\l-s)(\l+s)}{(2\l-1)(2\l+1)}}\;f_{\l-1,\emm},
\label{eq:Omega matrix}\\[6pt]
\big(\thpH f\big)_{\l\emm} &= -i\omega\,f_{\l\emm}.
\label{eq:thornp matrix}
\end{align}
\end{subequations}
The $\l=0$ entries degenerate and are given separately in
Appendix~\ref{app:held ops}.  Equation \eqref{eq:Omega matrix} is simply
$\Oh=-2ia\,\mathsf C$ written out, with $\mathsf C$ the $\cos\theta$ matrix
\eqref{eq:cos theta rec}; it is quoted here only so that all five operators
appear together.

At $a=0$ every operator in \eqref{eq:held matrices} collapses to its diagonal
and the multiplication operators \eqref{eq:mult matrices} vanish identically,
so the Schwarzschild construction is completely decoupled in $\l$ --- which
is why the closed-form results of Sec.~\ref{sec:schw} exist at all.

\PZ{\textbf{Discrepancy in $\Oh$ --- please check.}  Equation
\eqref{eq:Omega matrix} above is what $\Oh=-2ia\cos\theta$ requires.  The
matrix elements coded as \texttt{$\Omega$H[a][f][l,m]} in the notebooks do
\emph{not} agree with it.  Comparing entry by entry (any $s,\l,\emm$ gives
the same ratios):
\begin{center}
\begin{tabular}{lc}
\hline
entry & (notebook)/(Eq.~\eqref{eq:Omega matrix})\\
\hline
diagonal, $f_{\l\emm}$      & $-1$\\
upper, $f_{\l+1,\emm}$      & $-1/\sqrt2$\\
lower, $f_{\l-1,\emm}$      & $+1/\sqrt2$\\
\hline
\end{tabular}
\end{center}
No overall normalisation reconciles these, so this is not a convention
difference.  The cleanest diagnostic: $\Oh$ has spin weight zero, so it is
multiplication by a function and its matrix in an orthonormal basis must be
\emph{symmetric}.  The notebook's is \emph{antisymmetric} --- the
$f_{\l+1,\emm}$ and $f_{\l-1,\emm}$ coefficients have equal magnitude and
opposite sign --- which no multiplication operator can be.  It looks as
though $\Oh$ was implemented by copying the three-term template used for
$\edthH$ and $\tah$, which are genuinely derivative- and spin-raising-type and
do carry that relative minus.  Impact: $\Oh$ enters only \eqref{eq:zeta n Held}
and the $\tah,\tabh,\Oh$ terms of Appendix~\ref{app:Lie}, all of which vanish
at $a=0$ --- so the Schwarzschild results of Sec.~\ref{sec:schw} are
unaffected, but the Kerr gauge vector and $\Lie_{\zeta}g_{ab}$ are.}

\PZ{\textbf{And a related question about $\tah$.}  The notebook's $\tah$
carries the same opposite-sign pattern between its $f_{\l\pm1,\emm}$ entries.
For $\tah$ that is not immediately damning, since it maps the $\sYY{s}{}{}$
basis to the $\sYY{s+1}{}{}$ basis and the relevant condition is a transpose
relation against $\tabh$ rather than plain symmetry.  But given the $\Oh$
problem it is worth checking explicitly, e.g.\ by verifying
$\langle \sYY{s+1}{\l}{\emm}, \tah\, \sYY{s}{\l'}{\emm}\rangle
= \overline{\langle \sYY{s}{\l'}{\emm}, \tabh\, \sYY{s+1}{\l}{\emm}\rangle}$
once you have the closed form of $\tah$ for \eqref{eq:bkgd held}.}

\subsection{Inverse powers of $\rho$ and the $\cos\theta$ coupling matrix}
\label{sec:cos theta}
% ROADMAP: this is the step that makes the Kerr case genuinely different
% from Schwarzschild, and it deserves its own subsection.  Say: in
% Schwarzschild 1/rho = -r is a number, so (3.x) is trivial; in Kerr
% 1/rho = -r + i a cos theta mixes neighbouring l.  Hence Phi^adj at
% output l needs input up to l+3.
In Kerr, $\rho^{-1} = -(r-ia\cos\theta)$ is not a multiple of the identity in
$\l$-space.  Using
\begin{equation}\label{eq:cos theta rec}
\cos\theta\;\sYY{s}{\l}{\emm}
=
c_{\l+1}\,\sYY{s}{\l+1}{\emm}
+b_{\l}\,\sYY{s}{\l}{\emm}
+c_{\l}\,\sYY{s}{\l-1}{\emm},
\end{equation}
with
\begin{equation}\label{eq:cos theta coeffs}
c_{\l}(s,\emm)=\sqrt{\frac{(\l^{2}-\emm^{2})(\l^{2}-s^{2})}{\l^{2}(4\l^{2}-1)}},
\qquad
b_{\l}(s,\emm)=-\frac{s\,\emm}{\l(\l+1)},
\end{equation}
the operator $\rho^{-1}$ is represented by the tridiagonal matrix
\begin{equation}\label{eq:R matrix}
\mathsf R = -r\,\mathbb{1} + i a\,\mathsf C,
\qquad
\mathsf C_{L\l} = \text{matrix of \eqref{eq:cos theta rec}},
\end{equation}
and the adjusted Hertz potential \eqref{eq:Phi adj} is, as a vector in $\l$
at fixed $(\emm,n,k)$,
\begin{equation}\label{eq:Phi adj vector}
\boldsymbol{\Phi}^{\rm adj}
=
\boldsymbol{\Phi}^{0\circ}
+\mathsf R\,\boldsymbol{\Phi}^{1\circ}
+\mathsf R^{2}\,\boldsymbol{\Phi}^{2\circ}
+\mathsf R^{3}\,\boldsymbol{\Phi}^{3\circ}.
\end{equation}
Because $\mathsf R$ is tridiagonal, an output at multipole $L$ requires input
up to $L+3$; the implementation therefore pads the $\l$-range by three.
\PZ{The notebook pads by 3 for $\Phiadj$ and by \texttt{"LPad"}$\,=8$ for
$\zeta$.  Say why the $\zeta$ padding is larger (more nested operators), and
quote the observed truncation error for the padding you actually use.}

\subsection{Reality and conjugation in mode space}
\label{sec:reality}
% ---------------------------------------------------------------------
% ROADMAP.  Short but load-bearing: every sign error in an implementation
% of this scheme will come from here.  Three points: the conjugation rule,
% where the (-1)^s is not 1, and the practical shortcut of computing
% conjugate quantities from their own defining equation rather than by
% conjugating a mode array.
% ---------------------------------------------------------------------
Conjugation is the single most error-prone step in an implementation of
Sec.~\ref{sec:adjusted}, and we restate the rule here in the form in which it
is applied.  For a field of spin weight $s$ with $\sYY{s}{\l}{\emm}$
coefficients $X_{\lmnk}$, the conjugate field has spin weight $-s$ and
$\sYY{-s}{\l}{\emm}$ coefficients
\begin{equation}\label{eq:conj mode modes}
\big(\bar X\big)_{\lmnk}
=
(-1)^{s+\emm}\;\overline{X_{\l,-\emm,-n,-k}} ,
\end{equation}
the reflection $(\emm,n,k)\to(-\emm,-n,-k)$ supplying the frequency flip
$\wmnk\to-\wmnk$.  Thus $\bar\Phi$ at mode $(\lmnk)$ is built from $\Phi$ at
$(\l,-\emm,-n,-k)$, and similarly for $\fbHo$.

Two cautions follow.  First, the sign $(-1)^{s}$ is unity for all of
$\psifive$, $\fHo$ and $\PhiN{n}$, so it can be omitted throughout the
Hertz-potential sector without consequence; it is \emph{not} unity for
$\zetam$, for the intermediate chains $\edthH\fbHo$ and
$\edthpH\edthH^{2}\fbHo$, or for the $nm$ and $n\bar m$ metric components
(Table~\ref{tab:spin weights}).  Second, where a conjugate quantity satisfies
a defining equation of its own it is cheaper and safer to use it: we obtain
$\zetam$ from
\begin{equation}\label{eq:zeta m direct}
\zetam = \edthpH\fHo,
\end{equation}
the GHP conjugate of \eqref{eq:zeta m Held}, rather than by applying
\eqref{eq:conj mode modes} to $\zetamb$.

\PZ{\textbf{Please check two things in the notebooks.}
(a) \texttt{BondiGauge\_Adjusted.nb} defines
\texttt{cconj[expr] := ComplexExpand[Conjugate[expr]] /. \{$\omega\to-\omega$\}},
i.e.\ the frequency flip is inside \texttt{cconj}; the generic notebook
defines \texttt{cconj} \emph{without} the flip (it has a separate
\texttt{cconj\$$\omega$} that includes it) and carries the flip instead through
the $(-\emm,-n,-k)$ argument reflection.  Both are self-consistent, but
\texttt{cconjSch} differs between the two notebooks in exactly this way,
and any expression copied from one into the other will be wrong by
$\omega\to-\omega$.
(b) The generic notebook defines both \texttt{cconjMode} (with $(-1)^{\emm}$)
and \texttt{cconjMode\$s} (with $(-1)^{\emm+s}$).  Only the latter is
correct in general; confirm that every odd-spin-weight use calls the
\texttt{\$s} variant, or that the missing $(-1)^{s}$ has been inserted by
hand at the call site, as it appears to have been in
\texttt{Re$\S$dagMode$\Phi$Ops["nm"]} and \texttt{LieZetagComps["nmb"]}.}


% =====================================================================
\section{Schwarzschild implementation}
\label{sec:schw}
% =====================================================================
% ---------------------------------------------------------------------
% ROADMAP.  Everything is closed form here, so this section should read as
% a worked example: give the explicit formulas, then show the cancellation,
% then the numerics.  This is also where the reader can check the whole
% paper by hand, so do not compress it.
% ---------------------------------------------------------------------

\subsection{Closed-form mode formulas}
\label{sec:schw closed form}
In Schwarzschild the spheroidal harmonics reduce to spherical harmonics and
the mixing coefficients in \eqref{eq:spheroidal to spherical} become
$\delta_{L\l}$.  The angular constants reduce to
\begin{equation}\label{eq:B schw}
\Bco{2} = \frac{1}{4}\frac{(\l+2)!}{(\l-2)!}
= \frac{1}{4}(\l+2)(\l+1)\l(\l-1) = \frac{\LamL}{4},
\qquad
\Bco{1} = \frac{1}{2}\frac{(\l+1)!}{(\l-1)!} = \frac{\l(\l+1)}{2},
\end{equation}
with $\LamL\defeq(\l-1)\l(\l+1)(\l+2)$, and the denominator
\eqref{eq:Dlw def} becomes
\begin{equation}\label{eq:D schw}
\Dlw \;\longrightarrow\; \Bco{2}^{2}+9M^{2}\omega^{2}
= \frac{\LamL^{2}}{16}+9M^{2}\omega^{2},
\end{equation}
which we continue to denote $\Dlw$.  The seed and the four Held integration
constants are then
\begin{subequations}\label{eq:schw closed form}
\begin{align}
\fHo_{\l\emm}(\omega)
&=
\frac{2i\omega^{3}}{\Dlw}\;\psifiveM{\l\emm}(\omega),
\label{eq:f schw}\\[4pt]
\PhiN{0}_{\l\emm}(\omega)
&=
\frac{-\tfrac23 i\omega(-1)^{\emm}}{\Dlw}
\Big[\Bco{2}\Bco{1}-3iM\omega\big(1+\sqrt{\Bco{2}}\big)\Big]
\bar\psi^{5\circ}_{0,\l,-\emm}(-\omega),
\label{eq:Phi0 schw}\\[4pt]
\PhiN{1}_{\l\emm}(\omega)
&=
\frac{2\omega^{2}(-1)^{\emm}}{\Dlw}
\big[\Bco{2}-3iM\omega\big]
\bar\psi^{5\circ}_{0,\l,-\emm}(-\omega),
\label{eq:Phi1 schw}\\[4pt]
\PhiN{2}_{\l\emm}(\omega)
&=
\frac{4i\omega^{3}\,\Bco{2}/\Bco{1}\,(-1)^{\emm}}{\Dlw}
\bar\psi^{5\circ}_{0,\l,-\emm}(-\omega),
\label{eq:Phi2 schw}\\[4pt]
\PhiN{3}_{\l\emm}(\omega)
&=
\frac{-4\omega^{4}(-1)^{\emm}}{\Dlw}
\bar\psi^{5\circ}_{0,\l,-\emm}(-\omega),
\label{eq:Phi3 schw}
\end{align}
\end{subequations}
and the residual gauge vector is
\begin{subequations}\label{eq:zeta schw}
\begin{align}
\zetamM{\l\emm}(\omega)
&=
\frac{-2i\omega^{3}\sqrt{\Bco{2}/\Bco{1}}}{\Dlw}
\;\psifiveM{\l\emm}(\omega),
\label{eq:zeta m schw}\\[4pt]
\zetalM{\l\emm}(\omega)
&=
\frac{\omega^{2}\sqrt{\Bco{2}}}{\Dlw}
\;\psifiveM{\l\emm}(\omega)\;+\;{\rm c.c.},
\label{eq:zeta l schw}\\[4pt]
\zetanM{\l\emm}(\omega)
&=
\frac{-\omega^{2}\sqrt{\Bco{2}}}{\Dlw}
\Big(\Bco{1}+\tfrac12\Big)
\;\psifiveM{\l\emm}(\omega)\;+\;{\rm c.c.},
\label{eq:zeta n schw}
\end{align}
\end{subequations}
where ``$+\,{\rm c.c.}$'' denotes the mode-space real part \eqref{eq:Re mode},
i.e.\ the addition of $(-1)^{\emm}\bar\psi^{5\circ}_{0,\l,-\emm}(-\omega)$
inside the same prefactor.
\PZ{\textbf{Discrepancy to resolve.}  Equation \eqref{eq:Phi0 schw} uses
$\Bco{2}\Bco{1} = \tfrac18(\l+2)(\l+1)^{2}\l^{2}(\l-1)$, following
Vahid's notes Eq.~(29).  \texttt{BondiGauge\_Adjusted.nb} instead codes
$\tfrac18(\l+2)(\l+1)\l^{2}(\l-1)^{2}$, i.e.\ $(\l+1)^2(\l-1)$ replaced by
$(\l+1)(\l-1)^2$ -- these differ for every $\l\ge2$.  Everything else in
\eqref{eq:schw closed form}--\eqref{eq:zeta schw} agrees with the notebook up
to the uniform factor of 2 in the $\PhiN{n}$ noted in
Sec.~\ref{sec:conventions}.  The cleanest test is whether the $O(r)$
cancellation in Sec.~\ref{sec:schw verification} still holds with the notes'
coefficient.}

Note that $\sqrt{\Bco{2}}=\tfrac12\sqrt{\LamL}$ and
$\Bco{2}/\Bco{1} = \tfrac12(\l+2)(\l-1)$, which is the form in which these
factors appear in the implementation.

\subsection{Reconstruction and Lie-derivative operators in modes}
\label{sec:schw operators}
% ROADMAP: give the S^dagger components (App. C) and the L_zeta g components
% (App. D) specialised to a = 0, and say in one sentence that both are
% now ordinary radial ODEs in the single variable r with the l-dependence
% appearing only through Lambda_l.
The explicit Schwarzschild mode expressions for
$(\S^{\dagger}\Phi)_{ab}$ and $(\Lie_{\zeta}g)_{ab}$ are collected in
Appendices~\ref{app:Sdag} and \ref{app:Lie}.  Both are diagonal in $\l$ and
depend on it only through $\LamL$.

\subsection{Verification of the falloff}
\label{sec:schw verification}
% ROADMAP: this is the "proof" the paper offers.  Do it symbolically:
% expand h^adj_{ab} in powers of r and show the O(r), O(r^0) coefficients
% vanish.  Present as a table:  component | O(r) coeff | O(r^0) coeff |
% leading surviving power.  Include at least one nonaxisymmetric mode and
% one m = 0 mode, since the Re-projection differs.
\PZ{Table here.  Columns: NP component; coefficient of $r$; coefficient of
$r^{0}$; leading surviving power.  Rows: $nn$, $nm$, $n\bar m$, $mm$,
$\bar m\bar m$, for say $(\l,\emm)=(2,2),(2,0),(3,2),(3,-3)$.  Your
notebook already computes exactly these via
\texttt{Coefficient[hAdjustment[comp][l,m][r],r]}.}

An independent check is provided by the ``old'' and ``new'' implementations
of $\S^{\dagger}$ in the accompanying notebooks: one applies the operator
directly in closed form, the other builds it from the Held operators of
Sec.~\ref{sec:held operators} and reduces at the end.  The two agree to
machine precision for all modes tested.
\PZ{Quote the actual agreement level (the notebook does
\texttt{SetPrecision[Chop[checkTable],32]}).  This is a good, cheap
validation to report.}

\subsection{Results for circular equatorial orbits}
\label{sec:schw results}
\PZ{TODO -- numerics.  Suggested content: (i) a plot of the NP components of
$h^{\rm ret}$ and $h^{\rm adj}$ versus $r$ out to $r\sim10^{3}M$, showing the
$O(r)$ growth removed; (ii) a convergence plot in $\l_{\max}$; (iii) the
orbital parameters used ($r_{0}=10M$ in the notebook) and the source of the
$\psifive$ data (\texttt{psi0LargeRdata\_lmax40\_rOutbd1212.wxf}).
Please also state how $\psifive$ was extracted -- a large-$r$ fit to the
radial solution, or the $\mathscr I$ amplitude directly?  The notebook
filename suggests a fit, which is worth documenting since it sets the
achievable precision.}


% =====================================================================
\section{Kerr implementation}
\label{sec:kerr}
% =====================================================================
% ---------------------------------------------------------------------
% ROADMAP.  The structure mirrors Sec. 5 but the emphasis shifts: nothing
% is closed form, so the section is about the *pipeline* and its cost.
% Make the cost statement explicit -- it is the practical selling point:
%   per (m,n,k): one spheroidal expansion per l, then O(l_max) tridiagonal
%   applications.  Negligible next to the Teukolsky solves.
% ---------------------------------------------------------------------

\subsection{What changes relative to Schwarzschild}
\label{sec:kerr pipeline}
% ROADMAP: open by naming the three sources of l-mixing.  A reader who has
% followed Sec. V needs to know exactly where Kerr stops being diagonal,
% and there are precisely three places -- no more.
In Schwarzschild every step of Sec.~\ref{sec:algorithm} is diagonal in $\l$.
In Kerr, exactly three ingredients break that diagonality, and they are the
whole of the additional work:
\begin{enumerate}[label=(\alph*),leftmargin=*]
\item the input Teukolsky solution is expanded in spin-weighted
\emph{spheroidal} harmonics, whose overlap with the spherical basis spreads
each $\l$ over a band of $L$ (Sec.~\ref{sec:kerr mixing});
\item the Held angular operators $\edthH$, $\edthpH$ acquire off-diagonal
$a\omega$ terms, and the multiplication operators $\tah$, $\tabh$, $\Oh$ ---
absent in Schwarzschild --- are themselves tridiagonal
(Sec.~\ref{sec:kerr chains});
\item $\rho^{-1}=-(r-ia\cos\theta)$ is no longer proportional to the
identity, so assembling $\Phiadj$ from the $\PhiN{n}$ mixes three further
neighbours (Sec.~\ref{sec:kerr assembly}).
\end{enumerate}
Everything else --- the algebraic inversion for $\fHo$, the triangular system
for the $\PhiN{n}$, the $\thpH\to-i\omega$ replacement --- is unchanged.  In
particular every operator remains \emph{diagonal in $\emm$}, so the
completion commutes with an $\emm$-mode decomposition of the field; the
$\emm$-mode scheme of Ref.~\cite{Dolan:2012jg} may be used without
modification.

\subsection{Asymptotic amplitudes for generic bound orbits}
\label{sec:generic orbits}
% ROADMAP: the mode lattice, the canonical half, the caching structure.
% This is the expensive part of the whole calculation and deserves an
% honest cost statement.
For a generic bound geodesic the frequencies are \eqref{eq:omega mnk} and the
amplitudes $\Zinf{+2}{\lmnk}$ are computed mode by mode over the lattice
\begin{equation}\label{eq:mode lattice}
|s|\le\l\le\l_{\max},
\qquad
|\emm|\le\l,
\qquad
|n|\le n_{\max},
\qquad
|k|\le k_{\max}.
\end{equation}
Only the canonical half of the lattice,
\begin{equation}\label{eq:canonical half}
\emm>0
\quad\text{or}\quad
\big(\emm=0 \ \text{and}\ [\,n>0 \ \text{or}\ (n=0\ \text{and}\ k\ge0)\,]\big),
\end{equation}
need be computed directly; the remaining amplitudes follow from the
reflection symmetry \eqref{eq:Z symmetry}.  This halves the dominant cost of
the calculation, which is the set of Teukolsky radial solves and not the
completion itself (Sec.~\ref{sec:kerr cost}).
\PZ{State the $(n,k)$ truncation used in production and how it was chosen
(residual amplitude threshold?).  The notebook's exploratory values are
$\l_{\max}=4$, $n_{\max}=k_{\max}=4$, $p=10M$, $e=0.1$,
$\cos\chi=\cos(\pi/4)$, $a=0.667M$.  Worth also describing the per-$\l$
caching: the implementation writes one file per $\l$ as it goes and merges
into a global amplitude table, so a run is restartable --- a practical point
that belongs in the text if you want others to reproduce this.}

\subsection{Spheroidal-to-spherical mixing}
\label{sec:kerr mixing}
% ROADMAP: this is ingredient (a).  Give the projection, the band structure,
% and the truncation.  The key practical statement is that the b's are
% computed ONCE per (s, l, m, a omega) and memoised -- every Held chain
% downstream reuses them.
The seed $\fHo$ of \eqref{eq:f mode kerr} is obtained in the spheroidal
basis, one amplitude per $(\lmnk)$, exactly as in Schwarzschild.  It is then
projected onto the spherical basis using \eqref{eq:spheroidal to spherical}:
writing $\fHo_{\lmnk}$ for the spheroidal amplitude, the spherical
coefficients are
\begin{equation}\label{eq:f Y basis}
\fHo_{L\emm nk}
=
\sum_{\l}\fHo_{\lmnk}\;b^{(+2,\l,\emm)}_{L}(a\wmnk),
\end{equation}
and likewise for $\fbHo$ with $s=-2$.  Because the mixing coefficients depend
on $(s,\l,\emm,a\omega)$ only, they are computed once per mode and reused by
every Held chain downstream; in practice this memoization dominates the
efficiency of the whole completion, since the chains of
Sec.~\ref{sec:kerr chains} would otherwise recompute the same expansion
$O(10)$ times per mode.

Denote by $N_{\rm mix}$ the half-width at which the expansion is truncated,
\begin{equation}\label{eq:mixing truncation}
b^{(s,\l,\emm)}_{L}\ \text{retained for}\ |L-\l|\le N_{\rm mix}.
\end{equation}
\PZ{Quote $N_{\rm mix}$ and justify it.  The notebook uses
\texttt{"NumTerms"}$\,\to30$ with a chop tolerance of $10^{-25}$, which is
far more than needed for $a\omega\lesssim1$; a plot of $|b_{L}|$ against
$|L-\l|$ at the largest $a\wmnk$ occurring in your run would let you quote a
much smaller effective $N_{\rm mix}$ and would answer the truncation
question cleanly.  This also feeds directly into the padding budget of
Table~\ref{tab:bandwidth}.}

\subsection{Held chains as banded operators}
\label{sec:kerr chains}
% ROADMAP: ingredient (b).  The content here is (i) bandwidth accounting and
% (ii) the memoisation point.  Both are genuinely useful and neither is in
% Hollands-Toomani.
Each of $\edthH$, $\edthpH$, $\tah$, $\tabh$ and $\Oh$ is tridiagonal in $\l$
at fixed $\emm$ by \eqref{eq:held matrices}--\eqref{eq:mult matrices}, and
$\thpH$ is diagonal by \eqref{eq:thornp matrix}.  A chain of $p$ such
operators is therefore banded with half-bandwidth $p$: its value at output
multipole $L$ depends on the input at $L-p,\dots,L+p$.  Applying this to
\eqref{eq:Phi n Held} and \eqref{eq:zeta Held} gives the padding budget of
Table~\ref{tab:bandwidth}.

\begin{table}[htb]
\begin{ruledtabular}
\begin{tabular}{llc}
quantity & longest operator chain & half-bandwidth \\
\hline
$\PhiN{3}$ & $\thpH$ & $0$ \\
$\PhiN{2}$ & $\edthpH\edthH$ & $2$ \\
$\PhiN{1}$ & $\edthpH^{2}\edthH^{2}$ & $4$ \\
$\PhiN{0}$ & $\edthpH^{3}\edthH^{3}$ & $6$ \\
$\zetamb$  & $\edthH$ & $1$ \\
$\zetal$   & $\edthH^{2}$ & $2$ \\
$\zetan$   & $\edthpH\edthH^{3}$ & $4$ \\
$\rho^{-n}$ & $\mathsf C^{n}$ & $n\le3$ \\
\hline
$\Phiadj$ & $\max\{6,\,1{+}4,\,2{+}2,\,3{+}0\}$ & $6$ \\
\end{tabular}
\end{ruledtabular}
\caption{Half-bandwidth in $\l$ of each stage of the completion, at fixed
$(\emm,n,k)$.  To obtain $\Phiadj$ correctly for $L\le L_{\rm out}$ the seed
$\fbHo$ must be supplied for $\l\le L_{\rm out}+6$, and that seed must itself
be assembled from spheroidal amplitudes out to
$\l\le L_{\rm out}+6+N_{\rm mix}$.}
\label{tab:bandwidth}
\end{table}

Two implementation points follow.  First, because every operator couples
$L\to\{L-1,L,L+1\}$, a naive recursive evaluation of a chain of length $p$
costs $3^{p}$ leaf evaluations --- $729$ for $\PhiN{0}$ --- whereas
evaluating and caching each intermediate chain once, level by level, reduces
the cost to $O(p\,L_{\max})$.  We therefore build and store the intermediate
quantities
\begin{equation}\label{eq:chains}
\edthH\fbHo,\quad
\edthH^{2}\fbHo,\quad
\edthH^{3}\fbHo,\quad
\edthpH\edthH\fbHo,\quad
\edthpH\edthH^{2}\fbHo,\quad
\edthpH^{2}\edthH^{2}\fbHo,\quad
\edthpH^{3}\edthH^{3}\fbHo,
\end{equation}
together with $\tabh\,\edthH\fbHo$ and $\tabh\,\edthpH\edthH^{2}\fbHo$, and
assemble the $\PhiN{n}$ and $\zeta^{\circ}$ components from them.  Second,
the spin weight changes along each chain
($-2\to-1\to0\to+1$ under successive $\edthH$), and the matrix elements
\eqref{eq:held matrices} depend on it, so the spin weight must be tracked
with the intermediate rather than inferred at the end.

\subsection{Assembling $\Phiadj$}
\label{sec:kerr assembly}
% ROADMAP: ingredient (c).  Short -- the equation is already in Sec. IV D.
% What is worth adding here is the practical ordering (Horner) and the
% statement that this is where the +3 padding is consumed.
With the $\PhiN{n}$ in hand as vectors over $\l$ at fixed $(\emm,n,k)$, the
adjusted Hertz potential follows from \eqref{eq:Phi adj vector}.  It is
evaluated in Horner form,
\begin{equation}\label{eq:Phi adj horner}
\boldsymbol{\Phi}^{\rm adj}
=
\boldsymbol{\Phi}^{0\circ}
+\mathsf R\Big(\boldsymbol{\Phi}^{1\circ}
+\mathsf R\big(\boldsymbol{\Phi}^{2\circ}
+\mathsf R\,\boldsymbol{\Phi}^{3\circ}\big)\Big),
\end{equation}
which requires three tridiagonal multiplications rather than six and consumes
three multipoles of padding.  Note that $\mathsf R$ carries the explicit
radial dependence through $-r\,\mathbb{1}$, so \eqref{eq:Phi adj horner} is a
polynomial in $r$ of degree three with $\l$-banded matrix coefficients;
$\Phiadj$ and its first two radial derivatives, which are what
Appendix~\ref{app:Sdag} requires, are obtained from it analytically.

\subsection{The gauge vector in the spherical basis}
\label{sec:kerr zeta}
The components of $\zeta^{a}$ are built from the same cached chains.  Two
details differ from the Hertz sector.  First, $\zetal$ and $\zetan$ are real
spin-zero scalars and are assembled through the projection
\eqref{eq:Re mode}, which pairs the $(\emm,n,k)$ and $(-\emm,-n,-k)$ chains;
$\zetam$ is instead obtained directly from \eqref{eq:zeta m direct} to avoid
an explicit spin-weighted conjugation.  Second, the minimum output multipole
depends on the component: nearest-neighbour mixing can feed the spin-zero
components down to $L=|\emm|$, whereas the spin-$(\pm1)$ components start at
\begin{equation}\label{eq:Lmin}
L_{\min}
=
\begin{cases}
|\emm| , & \text{components } \zetal,\ \zetan \quad (s=0),\\[2pt]
\max(1,|\emm|), & \text{components } \zetam,\ \zetamb \quad (s=\pm1).
\end{cases}
\end{equation}
Finally, since $\fHo\propto\omega^{3}$, modes with $\wmnk=0$ contribute
nothing to the radiative completion and are skipped.
\PZ{Quote the $L$-padding actually used for $\zeta$ and the chop tolerance
below which coefficients are discarded.  Table~\ref{tab:bandwidth} suggests
$4$ suffices for $\zetan$; the notebook uses $8$, so either the extra margin
is unnecessary or I have miscounted the chain --- worth checking, since it
doubles the cost of this stage.}

\subsection{Cost and truncation}
\label{sec:kerr cost}
% ROADMAP: the selling point.  Make the comparison explicit and quantitative
% once you have timings.
At fixed $(\emm,n,k)$ the completion consists of one complex division per
$\l$ for the seed, one spheroidal projection per $\l$, and a fixed number
($\approx20$) of tridiagonal applications over the padded $\l$-range.  Its
cost is therefore $O(L_{\max})$ per $(\emm,n,k)$, against the
$O(L_{\max})$ \emph{Teukolsky radial solves} needed to produce the input
amplitudes in the first place, each of which is more expensive by orders of
magnitude.  The asymptotic completion is thus essentially free relative to
the calculation it corrects.
\PZ{Back this with one timing number --- e.g.\ wall-clock for the amplitude
run versus wall-clock for the completion at the same $\l_{\max}$.  A single
sentence with real numbers is worth a paragraph of argument here.}

\subsection{Results for circular equatorial orbits}
\label{sec:kerr results}
\PZ{TODO -- numerics.  Same three plots as Sec.~\ref{sec:schw results}, plus:
(i) a convergence study in $N_{\rm mix}$ at fixed $\l_{\max}$, which is the
one genuinely new truncation in Kerr; (ii) a spin sweep if you have it;
(iii) the $a\to0$ limit of the Kerr pipeline compared against the closed-form
Schwarzschild results of Sec.~\ref{sec:schw} --- that is the strongest single
validation available and costs nothing to run.}


% =====================================================================
\section{Discussion}
\label{sec:discussion}
% =====================================================================
% ---------------------------------------------------------------------
% ROADMAP.  Four short paragraphs:
%   1. What was done (one paragraph, no new content).
%   2. What it buys: the second-order source is now integrable at scri, so
%      asymptotic flux and memory formulas can be evaluated on the
%      reconstructed metric.  Connect to Hollands-Toomani Sec 3.3.2.
%   3. Limitations: stationary (omega = 0) completion still separate;
%      truncation in l for the cos-theta mixing; horizon-side analogue not
%      treated (the Sigma_- / H^- problem).
%   4. Outlook: second-order puncture, Lorenz-gauge transformation,
%      time-domain implementation.
% ---------------------------------------------------------------------
\PZ{TODO -- see the roadmap comment in the source.  Point 3 is the one to be
careful about: be explicit that this paper treats the $\scri$ side only, and
that the ingoing ($H^{-}$) counterpart is a separate construction.}


\begin{acknowledgments}
This work made use of the Black Hole Perturbation Toolkit~\cite{BHPToolkit},
in particular the \texttt{Teukolsky}, \texttt{KerrGeodesics} and
\texttt{SpinWeightedSpheroidalHarmonics} packages.
\PZ{Add funding acknowledgements.}
\end{acknowledgments}


% =====================================================================
\appendix
% =====================================================================

\section{Held angular operators: degenerate cases and conventions}
\label{app:held ops}
% ROADMAP: give the l = 0 entries omitted from (4.x), and state the
% convention for the square roots / branch choices.
For $\l=0$ the expressions \eqref{eq:held matrices}--\eqref{eq:mult matrices}
degenerate to
\begin{subequations}\label{eq:held l0}
\begin{align}
\big(\edthH f\big)_{0\emm} &= a\omega\sqrt{\tfrac{s(s-1)}{6}}\;f_{1\emm},
&
\big(\edthpH f\big)_{0\emm} &= a\omega\sqrt{\tfrac{s(s+1)}{6}}\;f_{1\emm},\\
\big(\tah f\big)_{0\emm} &= ia\sqrt{\tfrac{s(s-1)}{6}}\;f_{1\emm},
&
\big(\tabh f\big)_{0\emm} &= ia\sqrt{\tfrac{s(s+1)}{6}}\;f_{1\emm},\\
\big(\Oh f\big)_{00} &= -\frac{2ia}{\sqrt3}\;f_{10}.
&&
\end{align}
\end{subequations}
The $\edthH$, $\edthpH$, $\tah$ and $\tabh$ entries are transcribed from the
\texttt{If[l==0,...]} branches of the implementation.  The $\Oh$ entry is
instead the $\l=0$ limit of \eqref{eq:Omega matrix}: only $s=0$, $\emm=0$
occurs there, and $c_{1}(0,0)=1/\sqrt3$.
\PZ{The notebook's $\l=0$ branch for $\Oh$ reads
$ia\sqrt{-2(s-1)(s+1)/3}\,f_{1,0}$, which at $s=0$ gives $+ia\sqrt{2/3}$
rather than the $-2ia/\sqrt3$ above --- the same disagreement flagged after
\eqref{eq:Omega matrix}, carried into the degenerate case.  Since only $s=0$
can occur at $\l=0$, the $s$-dependence in that branch is inert in any event.
These entries matter only for the spin-zero gauge-vector components, where
Kerr nearest-neighbour mixing can feed $\l=0,1$; cf.\ the $L_{\min}$ switch
\eqref{eq:Lmin}.}

\section{Teukolsky--Starobinsky constants}
\label{app:TS}
For $s=\pm2$ the complex Teukolsky--Starobinsky constant is
\begin{equation}\label{eq:TS constant}
\Cts_{\l\emm}(a\omega) = D_{\l\emm}(a\omega) + 12\,i\,P\,\tilde\omega ,
\end{equation}
with
\begin{equation}\label{eq:D TS}
D_{\l\emm}
=
\Big[
\lambda_{\rm Ch}^{2}(\lambda_{\rm Ch}-2)^{2}
+8a\tilde\omega(\emm-a\tilde\omega)(\lambda_{\rm Ch}-2)(5\lambda_{\rm Ch}-4)
+48(a\tilde\omega)^{2}\big(2(\lambda_{\rm Ch}-2)+3(\emm-a\tilde\omega)^{2}\big)
\Big]^{1/2},
\end{equation}
where $\lambda_{\rm Ch}=\lambda+s^{2}+s$ and $\lambda$ is the spin-weighted
spheroidal eigenvalue, and $P=\pm1$ is the parity factor of
\eqref{eq:Zplus from Zminus}; see Ref.~\cite{Casals:2020fsb} for a careful
discussion of these constants and of the conventions attached to $P$.
In Schwarzschild $\lambda_{\rm Ch}=\l(\l+1)$,
$D_{\l}=\LamL$, and
\begin{equation}\label{eq:TS schw}
\Cts\,\bar\Cts = \LamL^{2}+144M^{2}\omega^{2}
= 16\big(\Bco{2}^{2}+9M^{2}\omega^{2}\big),
\end{equation}
which reproduces the Schwarzschild limit of \eqref{eq:f mode kerr} and fixes
the normalisation $\Cts\bar\Cts=16\,\Dlw$ used there.


\section{Reconstruction operator in modes (Schwarzschild)}
\label{app:Sdag}
% ROADMAP: give the IRG components of S^dagger acting on a single
% {}_{-2}Y_{lm} mode of Phi.  These are what the notebook codes.
With $\Phi=\sum_{\l\emm}\Phi_{\l\emm}(r)\,\sYY{-2}{\l}{\emm}$ in the ingoing
radiation gauge, the nonvanishing components of $\S^{\dagger}\Phi$ are
\begin{subequations}\label{eq:Sdag modes}
\begin{align}
\big(\S^{\dagger}\Phi\big)_{nn}
&= -\frac{1}{r^{2}}\sqrt{\frac{\LamL}{16}}
\Big[\Phi_{\l\emm}(r)+(-1)^{\emm}\bar\Phi_{\l,-\emm}(r)\Big],
\label{eq:Sdag nn}\\[4pt]
\big(\S^{\dagger}\Phi\big)_{nm}
&= -\frac{(-1)^{\emm}}{r^{2}}\sqrt{\frac{(\l-1)(\l+2)}{8}}\;
\overline{\Big[-2\Phi_{\l,-\emm}(r)+r\,\Phi'_{\l,-\emm}(r)\Big]},
\label{eq:Sdag nm}\\[4pt]
\big(\S^{\dagger}\Phi\big)_{n\bar m}
&= \frac{\sqrt{(\l-1)(\l+2)}}{2\sqrt2\,r^{2}}
\Big[-2\Phi_{\l\emm}(r)+r\,\Phi'_{\l\emm}(r)\Big],
\label{eq:Sdag nmb}\\[4pt]
\big(\S^{\dagger}\Phi\big)_{mm}
&= (-1)^{\emm}\,\overline{\Big[\frac{1}{r}\Phi'_{\l,-\emm}(r)
-\frac12\Phi''_{\l,-\emm}(r)\Big]},
\label{eq:Sdag mm}\\[4pt]
\big(\S^{\dagger}\Phi\big)_{\bar m\bar m}
&= \frac{1}{r}\Phi'_{\l\emm}(r)-\frac12\Phi''_{\l\emm}(r),
\label{eq:Sdag mbmb}
\end{align}
\end{subequations}
with
$\big(\S^{\dagger}\Phi\big)_{ll}=\big(\S^{\dagger}\Phi\big)_{ln}
=\big(\S^{\dagger}\Phi\big)_{lm}=\big(\S^{\dagger}\Phi\big)_{l\bar m}
=\big(\S^{\dagger}\Phi\big)_{m\bar m}=0$.
The signs multiplying the conjugated amplitudes in
\eqref{eq:Sdag nn}--\eqref{eq:Sdag mbmb} are not independent choices: each is
the factor $(-1)^{s+\emm}$ of the conjugation rule \eqref{eq:conj mode},
evaluated at the spin weight of the component being formed.  Thus
\eqref{eq:Sdag nn} and \eqref{eq:Sdag mm} carry $(-1)^{\emm}$, because
$h_{nn}$ has $s=0$ and $h_{mm}$ has $s=+2$, while \eqref{eq:Sdag nm} carries
$-(-1)^{\emm}$, because $h_{nm}$ has $s=+1$.
\PZ{Two things to check here.
(a) \emph{Overall factor.}  \eqref{eq:Sdag nn} already contains the
combination $\Phi+(-1)^{\emm}\bar\Phi$, which \emph{is} twice the mode-space
real part \eqref{eq:Re mode}, yet the notebook multiplies the whole set by a
further factor of 2.  Fix the convention once and propagate it; the falloff
cancellation is insensitive to an overall factor but the amplitudes are not.
(b) \emph{Which component carries the minus.}  In
\texttt{Re$\S$dagMode$\Phi$Ops} the $-(-1)^{\emm}$ sits on \texttt{"nm"} and
\texttt{"nmb"} is the primitive expression, whereas in
\texttt{LieZetagComps} it is the other way round --- \texttt{"nm"} is
primitive and \texttt{"nmb"} carries the sign.  Both are internally
consistent, but they can only be added together in \eqref{eq:h adj} if the
two conventions for which of $m^{a}$, $\bar m^{a}$ is contracted first agree.
Worth verifying explicitly, since an error here would flip the sign of one
$O(r)$ term and would not show up in the $nn$ or $mm$ channels.}

\section{Lie derivative of the metric along $\zeta^{a}$}
\label{app:Lie}
% ROADMAP: the long Held expressions.  These are transcribed from Vahid's
% notes eqs (18)-(20).  They are long enough that a transcription error is
% likely -- flag them and check against the notebook's LieZetagComps, which
% is the a = 0 specialisation.
With $\Lie_{\zeta}g_{ab}l^{a}=0$ and $\Lie_{\zeta}g_{ab}m^{a}\bar m^{b}=0$,
the remaining components are
\begin{align}
\Lie_{\zeta}g_{ab}\,m^{a}m^{b}
&=
-2\bar\rho\,\edthH\edthH\zetal
+2\edthpH\edthH\fHo
+2\Oh\thpH\fHo
+4\bar\rho\,\tah\,\edthpH\fHo
+8\rhopH\fHo ,
\label{eq:Lie mm}
\end{align}
\begin{align}
\Lie_{\zeta}g_{ab}\,n^{a}m^{b}
&=
\bar\rho\,\edthpH\edthH\edthH\zetal
-2\rho\bar\rho\,\tabh\,\edthH\edthH\zetal
+\rho(\rho+\bar\rho)\Psih\,\edthH\zetal
\non
&\quad
-\bar\rho\,\Oh\,\edthpH\edthpH\edthH\fHo
+2\rho\,\tabh\,\edthpH\edthH\fHo
\non
&\quad
+\frac{1}{\bar\rho}\Big(
-2\bar\rho\,\Oh{}^{2}\,\edthpH\thpH\fHo
+2(\rho+\bar\rho)\Oh\tabh\,\thpH\fHo
-2\bar\rho\,\tah\,\edthpH\edthpH\fHo
\non
&\qquad\quad
-(\rho^{2}\Oh+\rho)\Psih\,\edthpH\fHo
\Big)
\non
&\quad
+\bar\rho\,\Psibh\,\edthpH\fHo
-6\bar\rho\,\Oh\rhopH\,\edthpH\fHo
+4\rho\bar\rho\,\tah\tabh\,\edthpH\fHo
+8\rho\rhopH\tabh\,\fHo ,
\label{eq:Lie nm}
\end{align}
and
\begin{align}
\Lie_{\zeta}g_{ab}\,n^{a}n^{b}
&=
\Re\Big[
4\rho\bar\rho\,\tabh\,\edthpH\edthH\edthH\zetal
-4\rho^{2}\bar\rho\,\tabh{}^{2}\,\edthH\edthH\zetal
+2\rho^{2}\Psih\,\edthpH\edthH\zetal
+4\rho^{2}\bar\rho\,\Psih\tabh\,\edthH\zetal
\non
&\qquad
-2\edthpH\edthpH\edthpH\edthH\fHo
-\frac{2}{\rho}\Big(
4\Oh\,\edthpH\edthpH\thpH\fHo
-4\rho\bar\rho\,\Oh\tabh\,\edthpH\edthpH\edthH\fHo
\non
&\qquad\qquad
+4\big(4+\bar\rho\,\Oh-\rho\bar\rho\,\Oh{}^{2}\big)\tabh\,\edthpH\thpH\fHo
\Big)
\non
&\qquad
-\big(8\rhopH+3\rho\Psih+3\bar\rho\Psibh+2\rho^{2}\Oh{}^{2}\Psih
+8\rho\bar\rho\,\tah\tabh\big)\edthpH\edthpH\fHo
\non
&\qquad
+4\rho^{2}\tabh{}^{2}\,\edthpH\edthH\fHo
+4\rho(\rho+2\bar\rho)\Oh\tabh{}^{2}\,\thpH\fHo
+8\rho^{2}\bar\rho\,\tah\tabh{}^{2}\,\edthpH\fHo
\non
&\qquad
+4\rho\bar\rho\,\Psibh\tabh{}^{2}\,\edthpH\fHo
-24\rho\bar\rho\,\Oh\rhopH\tabh\,\edthpH\fHo
+16\rho^{2}\rhopH\tabh{}^{2}\,\fHo
\Big].
\label{eq:Lie nn}
\end{align}
The remaining components follow by conjugation,
$\Lie_{\zeta}g_{ab}n^{a}\bar m^{b}=\overline{\Lie_{\zeta}g_{ab}n^{a}m^{b}}$
and
$\Lie_{\zeta}g_{ab}\bar m^{a}\bar m^{b}=\overline{\Lie_{\zeta}g_{ab}m^{a}m^{b}}$,
which in mode space are \eqref{eq:conj mode} at $s=+1$ and $s=+2$
respectively, i.e.
\begin{equation}\label{eq:Lie conj modes}
\big(\Lie_{\zeta}g\big)_{n\bar m,\l\emm}
= -(-1)^{\emm}\,\overline{\big(\Lie_{\zeta}g\big)_{nm,\l,-\emm}},
\qquad
\big(\Lie_{\zeta}g\big)_{\bar m\bar m,\l\emm}
= (-1)^{\emm}\,\overline{\big(\Lie_{\zeta}g\big)_{mm,\l,-\emm}} .
\end{equation}
The relative minus sign between the two is the $(-1)^{s}$ of
\eqref{eq:conj mode} and is not optional.
\PZ{\textbf{Transcription warning.}  \eqref{eq:Lie mm}--\eqref{eq:Lie nn} are
transcribed from Eqs.~(18)--(20) of \texttt{metric-reconst.pdf}, where the
$n^{a}n^{b}$ expression in particular is typeset across a page break and the
grouping of the $1/\rho$ bracket in \eqref{eq:Lie nn} is ambiguous.  Please
check them against \texttt{LieZetagComps} in the notebook (which is the
$a\to0$ specialisation, so it will validate the $\Psih$, $\rhopH$ and
$\zetal$ terms but \emph{not} the $\tah$, $\tabh$, $\Oh$ terms).  A Kerr-side
check against Hollands--Toomani would be worth doing before submission.}

\PZ{The $a\to0$ specialisation actually used in the Schwarzschild numerics is
considerably shorter; consider giving it as a separate displayed set, since a
reader reproducing Sec.~\ref{sec:schw} needs only that.  I have not
transcribed it here to avoid duplicating an expression that is already long,
but it is the \texttt{"nn"/"nm"/"nmb"/"mm"/"mbmb"} entries of
\texttt{LieZetagComps} with $\rho=\bar\rho=-1/r$, $a=0$.}


\bibliographystyle{apsrev4-1}
\bibliography{MyReferences,extra}

\end{document}
